\input{preamble}

% OK, commencer ici.
%
\begin{document}

\title{Dualité pour les espaces}


\maketitle

\phantomsection
\label{section-phantom}

\tableofcontents

\section{Introduction}
\label{section-introduction}

\noindent
Ce chapitre est l'analogue du chapitre correspondant pour les
schémas, voir Dualité pour les schémas, section \ref{duality-section-introduction}.
Le développement est semblable à celui des articles
\cite{Neeman-Grothendieck}, \cite{LN},
\cite{Lipman-notes} et \cite{Neeman-improvement}.




\section{Complexes dualisants sur les espaces algébriques}
\label{section-dualizing-spaces}

\noindent
Soit $U$ un schéma localement noethérien. Soit $\mathcal{O}_\etale$
le faisceau structural de $U$ sur le petit site \'etale de $U$.
Nous dirons qu'un objet $K \in D_\QCoh(\mathcal{O}_\etale)$ est
un complexe dualisant sur $U$ si $K = \epsilon^*(\omega_U^\bullet)$
pour un complexe dualisant $\omega_U^\bullet$ au sens de
Dualité pour les schémas, section \ref{duality-section-dualizing-schemes}.
Ici, $\epsilon^* : D_\QCoh(\mathcal{O}_U) \to D_\QCoh(\mathcal{O}_\etale)$
est l'équivalence de Catégories dérivées des espaces, Lemme
\ref{spaces-perfect-lemma-derived-quasi-coherent-small-etale-site}.
La plupart des propriétés de $\omega_U^\bullet$ étudiées dans
Dualité pour les schémas, section \ref{duality-section-dualizing-schemes},
se transmettent à $K$ grâce à la discussion de
Catégories dérivées des espaces, sections
\ref{spaces-perfect-section-derived-quasi-coherent-etale} et
\ref{spaces-perfect-section-spell-out}.

\medskip\noindent
Nous définissons un complexe dualisant sur un espace algébrique localement noethérien
comme un complexe qui, localement pour la topologie \'etale, provient d'un complexe dualisant
sur le schéma correspondant.

\begin{lemma}
\label{lemma-equivalent-definitions}
Soit $S$ un schéma. Soit $X$ un espace algébrique localement noethérien sur $S$.
Soit $K$ un objet de $D_\QCoh(\mathcal{O}_X)$. Les conditions suivantes sont équivalentes :
\begin{enumerate}
\item Pour tout morphisme \'etale $U \to X$, où $U$ est un schéma,
la restriction $K|_U$ est un complexe dualisant sur $U$ (au sens expliqué ci-dessus).
\item Il existe un morphisme \'etale surjectif $U \to X$, où $U$ est un
schéma, tel que $K|_U$ soit un complexe dualisant sur $U$.
\end{enumerate}
\end{lemma}

\begin{proof}
Supposons que $U \to X$ soit \'etale surjectif, où $U$ est un schéma.
Soit $V \to X$ un morphisme \'etale, où $V$ est un schéma.
Alors
$$
U \leftarrow U \times_X V \rightarrow V
$$
sont des morphismes \'etales de schémas, celui qui aboutit à $V$ étant surjectif.
Nous pouvons donc appliquer Dualité pour les schémas, Lemme \ref{duality-lemma-descent-ascent},
pour voir que, si $K|_U$ est un complexe dualisant sur $U$, alors
$K|_V$ est un complexe dualisant sur $V$.
\end{proof}

\begin{definition}
\label{definition-dualizing-scheme}
Soit $S$ un schéma.
Soit $X$ un espace algébrique localement noethérien sur $S$.
Un objet $K$ de $D_\QCoh(\mathcal{O}_X)$ est appelé un
{\it complexe dualisant} si $K$ satisfait aux conditions équivalentes du
Lemme \ref{lemma-equivalent-definitions}.
\end{definition}

\begin{lemma}
\label{lemma-affine-duality}
Soit $A$ un anneau noethérien et soit $X = \Spec(A)$. Soit
$\mathcal{O}_\etale$ le faisceau structural de $X$ sur le
petit site \'etale de $X$. Soient $K, L$ des objets de $D(A)$.
Si $K \in D_{\textit{Coh}}(A)$ et si $L$ est de dimension injective
finie, alors
$$
\epsilon^*\widetilde{R\Hom_A(K, L)} =
R\SheafHom_{\mathcal{O}_\etale}(\epsilon^*\widetilde{K},
\epsilon^*\widetilde{L})
$$
dans $D(\mathcal{O}_\etale)$, où
$\epsilon : (X_\etale, \mathcal{O}_\etale) \to (X, \mathcal{O}_X)$
est comme dans Catégories dérivées des espaces, section
\ref{spaces-perfect-section-derived-quasi-coherent-etale}.
\end{lemma}

\begin{proof}
D'après Dualité pour les schémas, Lemme \ref{duality-lemma-affine-duality},
nous avons un isomorphisme canonique
$$
\widetilde{R\Hom_A(K, L)} =
R\SheafHom_{\mathcal{O}_X}(\widetilde{K}, \widetilde{L})
$$
dans $D(\mathcal{O}_X)$. Il existe un morphisme canonique
$$
\epsilon^*R\Hom_{\mathcal{O}_X}(\widetilde{K}, \widetilde{L})
\longrightarrow
R\SheafHom_{\mathcal{O}_\etale}(\epsilon^*\widetilde{K},
\epsilon^*\widetilde{L})
$$
dans $D(\mathcal{O}_\etale)$, voir Cohomologie sur les sites, Remarque
\ref{sites-cohomology-remark-prepare-fancy-base-change}.
Nous allons montrer que les membres de gauche et de droite de ce morphisme
ont des faisceaux de cohomologie isomorphes, mais nous omettrons de
vérifier que l'isomorphisme est donné par ce morphisme.

\medskip\noindent
Nous pouvons supposer que $L$ est donné par un complexe fini $I^\bullet$
de $A$-modules injectifs. Par récurrence sur la longueur de $I^\bullet$
et par compatibilité des constructions avec les triangles distingués,
nous nous ramenons au cas où $L = I[0]$, avec $I$ un $A$-module injectif.
Rappelons que les faisceaux de cohomologie de
$R\SheafHom_{\mathcal{O}_\etale}(\epsilon^*\widetilde{K},
\epsilon^*\widetilde{L}))$
sont les faisceaux associés au préfaisceau qui, à $U$ \'etale
sur $X$, associe le $i$-ème groupe d'Ext entre les restrictions de
$\epsilon^*\widetilde{K}$ et $\epsilon^*\widetilde{L}$
à $U_\etale$. Voir
Cohomologie sur les sites, Lemme
\ref{sites-cohomology-lemma-section-RHom-over-U}.
Si $U = \Spec(B)$ est affine, ce groupe d'Ext
est égal à $\text{Ext}^i_B(K \otimes_A B, L \otimes_A B)$
par l'équivalence de
Catégories dérivées des espaces, Lemme
\ref{spaces-perfect-lemma-derived-quasi-coherent-small-etale-site}, et
Catégories dérivées des schémas, Lemme
\ref{perfect-lemma-affine-compare-bounded}
(on utilise aussi les compatibilités détaillées dans
Catégories dérivées des espaces, Remarque
\ref{spaces-perfect-remark-match-total-direct-images}).
Puisque $A \to B$ est \'etale, nous voyons que
$I \otimes_A B$ est un $B$-module injectif
par Complexes dualisants, Lemme \ref{dualizing-lemma-injective-goes-up}.
Nous voyons donc que
\begin{align*}
\Ext^n_B(K \otimes_A B, I \otimes_A B)
& =
\Hom_B(H^{-n}(K \otimes_A B), I \otimes_A B) \\
& =
\Hom_{A_f}(H^{-n}(K) \otimes_A B, I \otimes_A B) \\
& =
\Hom_A(H^{-n}(K), I) \otimes_A B \\
& =
\text{Ext}^n_A(K, I) \otimes_A B
\end{align*}
L'avant-dernière égalité vaut parce que $H^{-n}(K)$ est un $A$-module fini, voir
Compléments d'algèbre, Lemme
\ref{more-algebra-lemma-pseudo-coherence-and-base-change-ext}.
Par conséquent, les faisceaux de cohomologie des membres de gauche et de droite
de l'égalité du lemme sont les mêmes.
\end{proof}

\begin{lemma}
\label{lemma-dualizing-spaces}
Soit $S$ un schéma. Soit $X$ un espace algébrique localement noethérien sur $S$.
Soit $K$ un complexe dualisant sur $X$.
Alors $K$ est un objet de $D_{\textit{Coh}}(\mathcal{O}_X)$
et $D = R\SheafHom_{\mathcal{O}_X}(-, K)$ induit une anti-équivalence
$$
D :
D_{\textit{Coh}}(\mathcal{O}_X)
\longrightarrow
D_{\textit{Coh}}(\mathcal{O}_X)
$$
munie d'un isomorphisme canonique
$\text{id} \to D \circ D$. Si $X$ est quasi-compact, alors
$D$ échange $D^+_{\textit{Coh}}(\mathcal{O}_X)$ et
$D^-_{\textit{Coh}}(\mathcal{O}_X)$ et induit une équivalence
$D^b_{\textit{Coh}}(\mathcal{O}_X) \to D^b_{\textit{Coh}}(\mathcal{O}_X)$.
\end{lemma}

\begin{proof}
Soit $U \to X$ un morphisme \'etale, avec $U$ affine. Écrivons $U = \Spec(A)$ et
soit $\omega_A^\bullet$ un complexe dualisant pour $A$ correspondant à $K|_U$
comme dans le Lemme \ref{lemma-equivalent-definitions} et
Dualité pour les schémas, Lemme \ref{duality-lemma-equivalent-definitions}.
D'après le Lemme \ref{lemma-affine-duality}, le diagramme
$$
\xymatrix{
D_{\textit{Coh}}(A) \ar[r] \ar[d]_{R\Hom_A(-, \omega_A^\bullet)} &
D_{\textit{Coh}}(\mathcal{O}_\etale)
\ar[d]^{R\SheafHom_{\mathcal{O}_\etale}(-, K|_U)} \\
D_{\textit{Coh}}(A) \ar[r] &
D(\mathcal{O}_\etale)
}
$$
est commutatif, où $\mathcal{O}_\etale$ est le faisceau structural du
petit site \'etale de $U$. Comme la formation de $R\SheafHom$ commute
à la restriction, nous concluons que $D$ envoie
$D_{\textit{Coh}}(\mathcal{O}_X)$ dans
$D_{\textit{Coh}}(\mathcal{O}_X)$. De plus, le morphisme canonique
$$
L \longrightarrow
R\SheafHom_{\mathcal{O}_X}(R\SheafHom_{\mathcal{O}_X}(L, K), K)
$$
(Cohomologie sur les sites, Lemme \ref{sites-cohomology-lemma-internal-hom-evaluate})
est un isomorphisme pour tout $L$ de $D_{\textit{Coh}}(\mathcal{O}_X)$,
car il en est ainsi au-dessus de tout $U$ comme ci-dessus, d'après
Complexes dualisants, Lemme \ref{dualizing-lemma-dualizing}.
L'assertion sur les propriétés de bornitude du foncteur $D$
dans le cas quasi-compact résulte également des assertions correspondantes de
Complexes dualisants, Lemme \ref{dualizing-lemma-dualizing}.
\end{proof}

\noindent
Soit $(\mathcal{C}, \mathcal{O})$ un site annelé. Rappelons
qu'un objet $L$ de $D(\mathcal{O})$ est {\it inversible}
s'il est un objet inversible pour la structure monoïdale symétrique
sur $D(\mathcal{O}_X)$ donnée par le produit tensoriel dérivé. Dans
Cohomologie sur les sites, Lemme \ref{sites-cohomology-lemma-invertible-derived},
nous avons vu que cela signifie que $L$ est parfait et que, si $(\mathcal{C}, \mathcal{O})$
est un site localement annelé, alors, pour tout objet $U$ de $\mathcal{C}$,
il existe un recouvrement $\{U_i \to U\}$ de $U$ dans $\mathcal{C}$
tel que $L|_{U_i} \cong \mathcal{O}_{U_i}[-n_i]$
pour certains entiers $n_i$.

\medskip\noindent
Soit $S$ un schéma et soit $X$ un espace algébrique sur $S$.
Si $L$ dans $D(\mathcal{O}_X)$ est inversible, il existe une
décomposition en réunion disjointe $X = \coprod_{n \in \mathbf{Z}} X_n$
telle que $L|_{X_n}$ soit un module inversible placé en degré $n$.
En particulier, il s'ensuit que $L = \bigoplus H^n(L)[-n]$,
ce qui donne un complexe bien défini de $\mathcal{O}_X$-modules
(à différentielles nulles) représentant $L$.

\begin{lemma}
\label{lemma-dualizing-unique-spaces}
Soit $S$ un schéma.
Soit $X$ un espace algébrique localement noethérien sur $S$.
Si $K$ et $K'$ sont des complexes dualisants sur $X$, alors $K'$
est isomorphe à $K \otimes_{\mathcal{O}_X}^\mathbf{L} L$
pour un objet inversible $L$ de $D(\mathcal{O}_X)$.
\end{lemma}

\begin{proof}
Posons
$$
L = R\SheafHom_{\mathcal{O}_X}(K, K')
$$
C'est un objet inversible de $D(\mathcal{O}_X)$, car cette assertion est vraie
localement sur les affines. Utiliser le Lemme \ref{lemma-affine-duality} et
Complexes dualisants, Lemme
\ref{dualizing-lemma-dualizing-unique}, ainsi que sa démonstration.
Le morphisme d'évaluation $L \otimes_{\mathcal{O}_X}^\mathbf{L} K \to K'$
est un isomorphisme pour la même raison.
\end{proof}

\begin{lemma}
\label{lemma-dimension-function-scheme}
Soit $S$ un schéma. Soit $X$ un espace algébrique localement noethérien
et quasi-séparé sur $S$.
Soit $\omega_X^\bullet$ un complexe dualisant sur $X$. Alors la fonction sur $X$
$|X| \to \mathbf{Z}$ définie par
$$
x \longmapsto \delta(x)\text{ tel que }
\omega_{X, \overline{x}}^\bullet[-\delta(x)]
\text{ soit un complexe dualisant normalisé sur }
\mathcal{O}_{X, \overline{x}}
$$
est une fonction de dimension sur $|X|$.
\end{lemma}

\begin{proof}
Soit $U$ un schéma et soit $U \to X$ un morphisme \'etale surjectif.
Soit $\omega_U^\bullet$ le complexe dualisant sur $U$ associé
à $\omega_X^\bullet|_U$.
Si $u \in U$ s'envoie sur $x \in |X|$, alors $\mathcal{O}_{X, \overline{x}}$
est l'hensélisé strict de $\mathcal{O}_{U, u}$. D'après
Complexes dualisants, Lemme \ref{dualizing-lemma-flat-unramified},
nous voyons que, si $\omega^\bullet$ est un complexe dualisant normalisé
pour $\mathcal{O}_{U, u}$, alors
$\omega^\bullet \otimes_{\mathcal{O}_{U, u}} \mathcal{O}_{X, \overline{x}}$
est un complexe dualisant normalisé pour $\mathcal{O}_{X, \overline{x}}$.
Nous voyons donc que la fonction de dimension $U \to \mathbf{Z}$ de
Dualité pour les schémas, Lemme \ref{duality-lemma-dimension-function-scheme},
pour le schéma $U$ et le complexe
$\omega_U^\bullet$ est égale au composé de $U \to |X|$ et de $\delta$.
En utilisant le fait que les spécialisations dans $|X|$ se relèvent en des spécialisations dans $U$
et que les spécialisations non triviales dans $U$ s'envoient sur des
spécialisations non triviales dans $X$
(Espaces raisonnables, Lemmes \ref{decent-spaces-lemma-decent-specialization} et
\ref{decent-spaces-lemma-decent-no-specializations-map-to-same-point}),
un argument topologique facile montre que $\delta$ est une fonction de dimension
sur $|X|$.
\end{proof}




 
\section{Adjoint à droite de l'image directe}
\label{section-twisted-inverse-image}

\noindent
C'est l'analogue de Dualité pour les schémas, section
\ref{duality-section-twisted-inverse-image}.

\begin{lemma}
\label{lemma-twisted-inverse-image}
\begin{reference}
C'est presque le même résultat que \cite[exemple 4.2]{Neeman-Grothendieck}.
\end{reference}
Soit $S$ un schéma.
Soit $f : X \to Y$ un morphisme entre espaces algébriques quasi-séparés et quasi-compacts
sur $S$. Le foncteur $Rf_* : D_\QCoh(X) \to D_\QCoh(Y)$
admet un adjoint à droite.
\end{lemma}

\begin{proof}
Nous allons démontrer l'existence d'un adjoint à droite en vérifiant les hypothèses de
Catégories dérivées, Proposition \ref{derived-proposition-brown}.
Tout d'abord, la catégorie $D_\QCoh(\mathcal{O}_X)$ admet les sommes directes, voir
Catégories dérivées des espaces, Lemme
\ref{spaces-perfect-lemma-quasi-coherence-direct-sums}.
La catégorie $D_\QCoh(\mathcal{O}_X)$ est compactement engendrée d'après
Catégories dérivées des espaces, Théorème
\ref{spaces-perfect-theorem-bondal-van-den-Bergh}.
Comme $X$ et $Y$ sont quasi-compacts et quasi-séparés, il en est de même de $f$, voir
Morphismes d'espaces, Lemmes
\ref{spaces-morphisms-lemma-compose-after-separated} et
\ref{spaces-morphisms-lemma-quasi-compact-permanence}.
Le foncteur $Rf_*$ commute donc aux sommes directes, voir
Catégories dérivées des espaces, Lemme
\ref{spaces-perfect-lemma-quasi-coherence-pushforward-direct-sums}.
Cela achève la démonstration.
\end{proof}

\begin{lemma}
\label{lemma-twisted-inverse-image-bounded-below}
Notations et hypothèses comme dans le Lemme \ref{lemma-twisted-inverse-image}.
Soit $a : D_\QCoh(\mathcal{O}_Y) \to D_\QCoh(\mathcal{O}_X)$ l'adjoint à droite
de $Rf_*$. Alors $a$ envoie
$D^+_\QCoh(\mathcal{O}_Y)$ dans $D^+_\QCoh(\mathcal{O}_X)$.
En fait, il existe un entier $N$ tel que
$H^i(K) = 0$ pour $i \leq c$ implique $H^i(a(K)) = 0$ pour $i \leq c - N$.
\end{lemma}

\begin{proof}
D'après Catégories dérivées des espaces, Lemme
\ref{spaces-perfect-lemma-quasi-coherence-direct-image},
le foncteur $Rf_*$ est de dimension cohomologique finie. Autrement dit,
il existe un entier $N$ tel que
$H^i(Rf_*L) = 0$ pour $i \geq N + c$ si $H^i(L) = 0$ pour $i \geq c$.
Soit $K \in D^+_\QCoh(\mathcal{O}_Y)$ tel que $H^i(K) = 0$ pour $i \leq c$.
Alors
$$
\Hom_{D(\mathcal{O}_X)}(\tau_{\leq c - N}a(K), a(K)) =
\Hom_{D(\mathcal{O}_Y)}(Rf_*\tau_{\leq c - N}a(K), K) = 0
$$
d'après ce qui précède. Cela implique manifestement que
$H^i(a(K)) = 0$ pour $i \leq c - N$.
\end{proof}

\noindent
Soit $S$ un schéma.
Soit $f : X \to Y$ un morphisme d'espaces algébriques quasi-séparés et quasi-compacts
sur $S$.
Notons $a$ l'adjoint à droite de
$Rf_* : D_\QCoh(\mathcal{O}_X) \to D_\QCoh(\mathcal{O}_Y)$. Pour tous
$K \in D_\QCoh(\mathcal{O}_Y)$ et $L \in D_\QCoh(\mathcal{O}_X)$,
nous obtenons un morphisme canonique
\begin{equation}
\label{equation-sheafy-trace}
Rf_*R\SheafHom_{\mathcal{O}_X}(L, a(K))
\longrightarrow
R\SheafHom_{\mathcal{O}_Y}(Rf_*L, K)
\end{equation}
Ce morphisme est en effet défini comme le composé
$$
Rf_*R\SheafHom_{\mathcal{O}_X}(L, a(K)) \to
R\SheafHom_{\mathcal{O}_Y}(Rf_*L, Rf_*a(K)) \to
R\SheafHom_{\mathcal{O}_Y}(Rf_*L, K)
$$
où la première flèche est celle de
Cohomologie sur les sites, Remarque
\ref{sites-cohomology-remark-projection-formula-for-internal-hom},
et la seconde est la co-unité $Rf_*a(K) \to K$ de l'adjonction.

\begin{lemma}
\label{lemma-iso-on-RSheafHom}
Soit $S$ un schéma.
Soit $f : X \to Y$ un morphisme d'espaces algébriques quasi-compacts et quasi-séparés
sur $S$.
Soit $a$ l'adjoint à droite de
$Rf_* : D_\QCoh(\mathcal{O}_X) \to D_\QCoh(\mathcal{O}_Y)$.
Soient $L \in D_\QCoh(\mathcal{O}_X)$ et
$K \in D_\QCoh(\mathcal{O}_Y)$.
Alors le morphisme (\ref{equation-sheafy-trace})
$$
Rf_*R\SheafHom_{\mathcal{O}_X}(L, a(K))
\longrightarrow
R\SheafHom_{\mathcal{O}_Y}(Rf_*L, K)
$$
devient un isomorphisme après application du foncteur
$DQ_Y : D(\mathcal{O}_Y) \to D_\QCoh(\mathcal{O}_Y)$
décrit dans Catégories dérivées des espaces, section
\ref{spaces-perfect-section-better-coherator}.
\end{lemma}

\begin{proof}
L'énoncé a un sens puisque $DQ_Y$ existe d'après
Catégories dérivées des espaces, Lemme
\ref{spaces-perfect-lemma-better-coherator}.
Comme $DQ_Y$ est l'adjoint à droite du foncteur d'inclusion
$D_\QCoh(\mathcal{O}_Y) \to D(\mathcal{O}_Y)$,
pour démontrer le lemme, il suffit de montrer que, pour tout
$M \in D_\QCoh(\mathcal{O}_Y)$,
le morphisme (\ref{equation-sheafy-trace}) induit une bijection
$$
\Hom_Y(M, Rf_*R\SheafHom_{\mathcal{O}_X}(L, a(K)))
\longrightarrow
\Hom_Y(M, R\SheafHom_{\mathcal{O}_Y}(Rf_*L, K))
$$
Pour le voir, nous utilisons la suite d'égalités suivante :
\begin{align*}
\Hom_Y(M, Rf_*R\SheafHom_{\mathcal{O}_X}(L, a(K)))
& =
\Hom_X(Lf^*M, R\SheafHom_{\mathcal{O}_X}(L, a(K))) \\
& =
\Hom_X(Lf^*M \otimes_{\mathcal{O}_X}^\mathbf{L} L, a(K)) \\
& =
\Hom_Y(Rf_*(Lf^*M \otimes_{\mathcal{O}_X}^\mathbf{L} L), K) \\
& =
\Hom_Y(M \otimes_{\mathcal{O}_Y}^\mathbf{L} Rf_*L, K) \\
& =
\Hom_Y(M, R\SheafHom_{\mathcal{O}_Y}(Rf_*L, K))
\end{align*}
La première égalité résulte de Cohomologie sur les sites, Lemme
\ref{sites-cohomology-lemma-adjoint}.
La deuxième résulte de Cohomologie sur les sites, Lemme
\ref{sites-cohomology-lemma-internal-hom}.
La troisième résulte de la construction de $a$.
La quatrième résulte de Catégories dérivées des espaces, Lemme
\ref{spaces-perfect-lemma-cohomology-base-change} (c'est l'étape importante).
La cinquième résulte de Cohomologie sur les sites, Lemme
\ref{sites-cohomology-lemma-internal-hom}.
\end{proof}

\begin{example}
\label{example-iso-on-RSheafHom}
L'énoncé du Lemme \ref{lemma-iso-on-RSheafHom} n'est pas vrai sans
appliquer le « cohérateur » $DQ_Y$. Voir
Dualité pour les schémas, Exemple \ref{duality-example-iso-on-RSheafHom}.
\end{example}

\begin{remark}
\label{remark-iso-on-RSheafHom}
Dans la situation du Lemme \ref{lemma-iso-on-RSheafHom}, nous avons
$$
DQ_Y(Rf_*R\SheafHom_{\mathcal{O}_X}(L, a(K))) =
Rf_* DQ_X(R\SheafHom_{\mathcal{O}_X}(L, a(K)))
$$
d'après Catégories dérivées des espaces, Lemme
\ref{spaces-perfect-lemma-pushforward-better-coherator}.
Ainsi, si $R\SheafHom_{\mathcal{O}_X}(L, a(K)) \in D_\QCoh(\mathcal{O}_X)$,
nous pouvons « supprimer » le $DQ_Y$ du membre de gauche du morphisme.
D'autre part, si nous savons que
$R\SheafHom_{\mathcal{O}_Y}(Rf_*L, K) \in D_\QCoh(\mathcal{O}_Y)$,
nous pouvons « supprimer » le $DQ_Y$ du membre de droite du morphisme.
Si ces deux conditions sont satisfaites, nous voyons que (\ref{equation-sheafy-trace})
est un isomorphisme. En combinant ceci avec
Catégories dérivées des espaces, Lemme
\ref{spaces-perfect-lemma-quasi-coherence-internal-hom},
nous voyons que $Rf_*R\SheafHom_{\mathcal{O}_X}(L, a(K)) \to
R\SheafHom_{\mathcal{O}_Y}(Rf_*L, K)$ est un isomorphisme si
\begin{enumerate}
\item $L$ et $Rf_*L$ sont parfaits, ou
\item $K$ est borné inférieurement et $L$ et $Rf_*L$ sont pseudo-cohérents.
\end{enumerate}
Pour (2), nous utilisons le fait que $a(K)$ est borné inférieurement si $K$
est borné inférieurement, voir Lemme \ref{lemma-twisted-inverse-image-bounded-below}.
\end{remark}

\begin{example}
\label{example-iso-on-RSheafHom-noetherian}
Soit $S$ un schéma.
Soit $f : X \to Y$ un morphisme propre d'espaces algébriques noethériens
sur $S$, $L \in D^-_{\textit{Coh}}(X)$ et $K \in D^+_{\QCoh}(\mathcal{O}_Y)$.
Alors le morphisme $Rf_*R\SheafHom_{\mathcal{O}_X}(L, a(K)) \to
R\SheafHom_{\mathcal{O}_Y}(Rf_*L, K)$ est un isomorphisme.
En effet, les complexes $L$ et $Rf_*L$ sont pseudo-cohérents d'après
Catégories dérivées des espaces, Lemmes
\ref{spaces-perfect-lemma-identify-pseudo-coherent-noetherian} et
\ref{spaces-perfect-lemma-direct-image-coherent},
et la discussion de la Remarque \ref{remark-iso-on-RSheafHom} s'applique.
\end{example}

\begin{lemma}
\label{lemma-iso-global-hom}
Soit $S$ un schéma.
Soit $f : X \to Y$ un morphisme d'espaces algébriques quasi-séparés et quasi-compacts
sur $S$.
Pour tous $L \in D_\QCoh(\mathcal{O}_X)$ et $K \in D_\QCoh(\mathcal{O}_Y)$,
(\ref{equation-sheafy-trace}) induit un isomorphisme
$R\Hom_X(L, a(K)) \to R\Hom_Y(Rf_*L, K)$ entre les Hom dérivés globaux.
\end{lemma}

\begin{proof}
Par construction (Cohomologie sur les sites, section
\ref{sites-cohomology-section-global-RHom}), les complexes
$$
R\Hom_X(L, a(K)) =
R\Gamma(X, R\SheafHom_{\mathcal{O}_X}(L, a(K))) =
R\Gamma(Y, Rf_*R\SheafHom_{\mathcal{O}_X}(L, a(K)))
$$
et
$$
R\Hom_Y(Rf_*L, K) = R\Gamma(Y, R\SheafHom_{\mathcal{O}_X}(Rf_*L, a(K)))
$$
Ainsi, le lemme est une conséquence du Lemme \ref{lemma-iso-on-RSheafHom}.
En effet, un morphisme $E \to E'$ dans $D(\mathcal{O}_Y)$ qui induit
un isomorphisme $DQ_Y(E) \to DQ_Y(E')$ induit un quasi-isomorphisme
$R\Gamma(Y, E) \to R\Gamma(Y, E')$. Nous avons en effet
$H^i(Y, E) = \Ext^i_Y(\mathcal{O}_Y, E) = \Hom(\mathcal{O}_Y[-i], E) =
\Hom(\mathcal{O}_Y[-i], DQ_Y(E))$, car $\mathcal{O}_Y[-i]$
appartient à $D_\QCoh(\mathcal{O}_Y)$ et $DQ_Y$ est l'adjoint à droite
du foncteur d'inclusion $D_\QCoh(\mathcal{O}_Y) \to D(\mathcal{O}_Y)$.
\end{proof}







 
\section{Adjoint à droite de l'image directe et changement de base, I}
\label{section-base-change-map}

\noindent
Définissons le morphisme de changement de base entre les adjoints à droite de l'image directe.
Soit $S$ un schéma. Considérons un diagramme cartésien
\begin{equation}
\label{equation-base-change}
\vcenter{
\xymatrix{
X' \ar[r]_{g'} \ar[d]_{f'} & X \ar[d]^f \\
Y' \ar[r]^g & Y
}
}
\end{equation}
où $Y'$ et $X$ sont {\bf indépendants au sens de Tor} sur $Y$. Notons
$$
a  : D_\QCoh(\mathcal{O}_Y) \to D_\QCoh(\mathcal{O}_X)
\quad\text{et}\quad
a' : D_\QCoh(\mathcal{O}_{Y'}) \to D_\QCoh(\mathcal{O}_{X'})
$$
les adjoints à droite de $Rf_*$ et $Rf'_*$
(Lemme \ref{lemma-twisted-inverse-image}).
Le morphisme de changement de base de
Cohomologie sur les sites, Remarque \ref{sites-cohomology-remark-base-change},
donne une transformation de foncteurs
$$
Lg^* \circ Rf_* \longrightarrow Rf'_* \circ L(g')^*
$$
sur les catégories dérivées de faisceaux à cohomologie quasi-cohérente.
Il donne donc une transformation entre les adjoints à droite en sens opposé
$$
a \circ Rg_* \longleftarrow Rg'_* \circ a'
$$

\begin{lemma}
\label{lemma-flat-precompose-pus}
Dans le diagramme (\ref{equation-base-change}), le morphisme
$a \circ Rg_* \leftarrow Rg'_* \circ a'$ est un isomorphisme.
\end{lemma}

\begin{proof}
Le morphisme de changement de base $Lg^* \circ Rf_* K \to Rf'_* \circ L(g')^*K$
est un isomorphisme pour tout $K$ de $D_\QCoh(\mathcal{O}_X)$ d'après
Catégories dérivées des espaces, Lemme
\ref{spaces-perfect-lemma-compare-base-change}
(on utilise ici l'hypothèse d'indépendance au sens de Tor).
La transformation correspondante entre foncteurs adjoints est donc
elle aussi un isomorphisme.
\end{proof}

\noindent
Nous pouvons alors considérer le
morphisme de foncteurs
$D_\QCoh(\mathcal{O}_Y) \to D_\QCoh(\mathcal{O}_{X'})$
donné par le composé
\begin{equation}
\label{equation-base-change-map}
L(g')^* \circ a \to
L(g')^* \circ a \circ Rg_* \circ Lg^* \leftarrow
L(g')^* \circ Rg'_* \circ a' \circ Lg^* \to a' \circ Lg^*
\end{equation}
La première flèche provient du morphisme d'adjonction $\text{id} \to Rg_* Lg^*$,
et la dernière du morphisme d'adjonction $L(g')^*Rg'_* \to \text{id}$.
Nous avons besoin de l'hypothèse d'indépendance au sens de Tor pour inverser la flèche
du milieu, voir Lemme \ref{lemma-flat-precompose-pus}.
On peut aussi considérer (\ref{equation-base-change-map}), par
adjonction entre $L(g')^*$ et $R(g')_*$, comme une transformation naturelle
$$
a \to a \circ Rg_* \circ Lg^* \leftarrow Rg'_* \circ a' \circ Lg^*
$$
où, là encore, la seconde flèche est inversible. Si $M \in D_\QCoh(\mathcal{O}_X)$
et $K \in D_\QCoh(\mathcal{O}_Y)$,
alors ce morphisme est donné sur les foncteurs de Yoneda par
\begin{align*}
\Hom_X(M, a(K))
& =
\Hom_Y(Rf_*M, K) \\
& \to
\Hom_Y(Rf_*M, Rg_* Lg^*K) \\
& =
\Hom_{Y'}(Lg^*Rf_*M, Lg^*K) \\
& \leftarrow
\Hom_{Y'}(Rf'_* L(g')^*M, Lg^*K) \\
& =
\Hom_{X'}(L(g')^*M, a'(Lg^*K)) \\
& =
\Hom_X(M, Rg'_*a'(Lg^*K))
\end{align*}
(la flèche dirigée vers la gauche est inversible par le théorème de
changement de base donné dans
Catégories dérivées des espaces, Lemme
\ref{spaces-perfect-lemma-compare-base-change}),
ce qui rend les choses un peu plus explicites.

\medskip\noindent
Dans cette section, nous démontrons d'abord que le morphisme de changement de base satisfait
à certaines compatibilités naturelles relativement à l'empilement de carrés, comme dans
Cohomologie sur les sites, Remarques
\ref{sites-cohomology-remark-compose-base-change} et
\ref{sites-cohomology-remark-compose-base-change-horizontal},
pour le morphisme de changement de base usuel.
Nous suggérons au lecteur de sauter le reste de cette section en première lecture.

\begin{lemma}
\label{lemma-compose-base-change-maps}
Soit $S$ un schéma. Considérons un diagramme commutatif
$$
\xymatrix{
X' \ar[r]_k \ar[d]_{f'} & X \ar[d]^f \\
Y' \ar[r]^l \ar[d]_{g'} & Y \ar[d]^g \\
Z' \ar[r]^m & Z
}
$$
d'espaces algébriques quasi-compacts et quasi-séparés sur $S$, dans lequel
les deux carrés sont cartésiens, et où $f$ et $l$
ainsi que $g$ et $m$ sont indépendants au sens de Tor.
Alors les morphismes (\ref{equation-base-change-map})
des deux carrés se composent pour donner le morphisme de
changement de base du rectangle extérieur (voir la démonstration pour un énoncé précis).
\end{lemma}

\begin{proof}
Il résulte des hypothèses que $g \circ f$ et $m$ sont indépendants au sens de Tor
(nous omettons les détails), de sorte que l'énoncé a un sens.
Dans cette démonstration, nous écrivons $k^*$ à la place de $Lk^*$ et $f_*$ à la place
de $Rf_*$. Soient $a$, $b$ et $c$ les adjoints à droite du
Lemme \ref{lemma-twisted-inverse-image}
pour $f$, $g$ et $g \circ f$, et de même pour les versions munies d'un prime.
La flèche correspondant au carré supérieur est le composé
$$
\gamma_{top} :
k^* \circ a \to k^* \circ a \circ l_* \circ l^*
\xleftarrow{\xi_{top}} k^* \circ k_* \circ a' \circ l^* \to a' \circ l^*
$$
où $\xi_{top} : k_* \circ a' \to a \circ l_*$
est un isomorphisme (et peut donc être inversé)
et est la flèche « duale » du morphisme de changement de base
$l^* \circ f_* \to f'_* \circ k^*$. Les flèches extérieures proviennent
des morphismes canoniques $1 \to l_* \circ l^*$ et $k^* \circ k_* \to 1$.
De même, pour le second carré, nous avons
$$
\gamma_{bot} :
l^* \circ b \to l^* \circ b \circ m_* \circ m^*
\xleftarrow{\xi_{bot}} l^* \circ l_* \circ b' \circ m^* \to b' \circ m^*
$$
Pour le rectangle extérieur, nous obtenons
$$
\gamma_{rect} :
k^* \circ c \to k^* \circ c \circ m_* \circ m^*
\xleftarrow{\xi_{rect}} k^* \circ k_* \circ c' \circ m^* \to c' \circ m^*
$$
Nous avons $(g \circ f)_* = g_* \circ f_*$ et donc
$c = a \circ b$, et de même $c' = a' \circ b'$.
L'assertion du lemme est que $\gamma_{rect}$
est égal au composé
$$
k^* \circ c = k^* \circ a \circ b \xrightarrow{\gamma_{top}}
a' \circ l^* \circ b \xrightarrow{\gamma_{bot}}
a' \circ b' \circ m^* = c' \circ m^*
$$
Pour le voir, considérons le diagramme suivant :
$$
\xymatrix{
& & k^* \circ a \circ b \ar[d] \ar[lldd] \\
& & k^* \circ a \circ l_* \circ l^* \circ b \ar[ld] \\
k^* \circ a \circ b \circ m_* \circ m^* \ar[r] &
k^* \circ a \circ l_* \circ l^* \circ b \circ m_* \circ m^* &
k^* \circ k_* \circ a' \circ l^* \circ b \ar[u]_{\xi_{top}} \ar[d] \ar[ld] \\
& k^*\circ k_* \circ a' \circ l^* \circ b \circ m_* \circ m^*
\ar[u]_{\xi_{top}} \ar[rd] &
a' \circ l^* \circ b \ar[d] \\
k^* \circ k_* \circ a' \circ b' \circ m^* \ar[uu]_{\xi_{rect}} \ar[ddrr] &
k^*\circ k_* \circ a' \circ l^* \circ l_* \circ b' \circ m^*
\ar[u]_{\xi_{bot}} \ar[l] \ar[dr] &
a' \circ l^* \circ b \circ m_* \circ m^* \\
& & a' \circ l^* \circ l_* \circ b' \circ m^* \ar[u]_{\xi_{bot}} \ar[d] \\
& & a' \circ b' \circ m^*
}
$$
En descendant du côté droit, nous obtenons le composé, et en descendant
du côté gauche, nous obtenons $\gamma_{rect}$.
Tous les quadrilatères du côté droit de ce diagramme sont commutatifs
d'après Catégories, Lemme \ref{categories-lemma-properties-2-cat-cats},
ou, plus simplement, d'après la discussion qui précède
Catégories, Définition \ref{categories-definition-horizontal-composition}.
Il suffit donc de montrer que le diagramme
$$
\xymatrix{
a \circ l_* \circ l^* \circ b \circ m_* &
a \circ b \circ m_* \ar[l] \\
k_* \circ a' \circ l^* \circ b \circ m_* \ar[u]_{\xi_{top}} & \\
k_* \circ a' \circ l^* \circ l_* \circ b' \ar[u]_{\xi_{bot}} \ar[r] &
k_* \circ a' \circ b' \ar[uu]_{\xi_{rect}}
}
$$
devient commutatif lorsque nous inversons les flèches $\xi_{top}$, $\xi_{bot}$
et $\xi_{rect}$ (noter que cela diffère du fait de demander que le
diagramme soit commutatif). Cependant, le diagramme
$$
\xymatrix{
& a \circ l_* \circ l^* \circ b \circ m_* \\
a \circ l_* \circ l^* \circ l_* \circ b'
\ar[ru]^{\xi_{bot}} & &
k_* \circ a' \circ l^* \circ b \circ m_* \ar[ul]_{\xi_{top}} \\
& k_* \circ a' \circ l^* \circ l_* \circ b'
\ar[ul]^{\xi_{top}} \ar[ur]_{\xi_{bot}}
}
$$
est commutatif d'après Catégories, Lemme \ref{categories-lemma-properties-2-cat-cats}.
Puisque les diagrammes
$$
\vcenter{
\xymatrix{
a \circ l_* \circ l^* \circ b \circ m_* & a \circ b \circ m \ar[l] \\
a \circ l_* \circ l^* \circ l_* \circ b' \ar[u] &
a \circ l_* \circ b' \ar[l] \ar[u]
}
}
\quad\text{et}\quad
\vcenter{
\xymatrix{
a \circ l_* \circ l^* \circ l_* \circ b' \ar[r] & a \circ l_* \circ b' \\
k_* \circ a' \circ l^* \circ l_* \circ b' \ar[u] \ar[r] &
k_* \circ a' \circ b' \ar[u]
}
}
$$
sont commutatifs (voir les références citées) et puisque le composé de
$l_* \to l_* \circ l^* \circ l_* \to l_*$ est l'identité,
nous trouvons qu'il suffit de démontrer que
$$
k \circ a' \circ b' \xrightarrow{\xi_{bot}} a \circ l_* \circ b
\xrightarrow{\xi_{top}} a \circ b \circ m_*
$$
est égal à $\xi_{rect}$ (via les identifications $a \circ b = c$
et $a' \circ b' = c'$). C'est l'assertion duale de
Cohomologie sur les sites, Remarque \ref{sites-cohomology-remark-compose-base-change},
ce qui achève la démonstration.
\end{proof}

\begin{lemma}
\label{lemma-compose-base-change-maps-horizontal}
Soit $S$ un schéma. Considérons un diagramme commutatif
$$
\xymatrix{
X'' \ar[r]_{g'} \ar[d]_{f''} & X' \ar[r]_g \ar[d]_{f'} & X \ar[d]^f \\
Y'' \ar[r]^{h'} & Y' \ar[r]^h & Y
}
$$
d'espaces algébriques quasi-compacts et quasi-séparés sur $S$, dans lequel
les deux carrés sont cartésiens, et où $f$ et $h$
ainsi que $f'$ et $h'$ sont indépendants au sens de Tor.
Alors les morphismes (\ref{equation-base-change-map})
des deux carrés se composent pour donner le morphisme de
changement de base du rectangle extérieur (voir la démonstration pour un énoncé précis).
\end{lemma}

\begin{proof}
Il résulte des hypothèses que $f$ et $h \circ h'$ sont indépendants au sens de Tor
(nous omettons les détails), de sorte que l'énoncé a un sens.
Dans cette démonstration, nous écrivons $g^*$ à la place de $Lg^*$ et $f_*$ à la place
de $Rf_*$. Soient $a$, $a'$ et $a''$ les adjoints à droite du
Lemme \ref{lemma-twisted-inverse-image}
pour $f$, $f'$ et $f''$. La flèche correspondant au carré de droite
est le composé
$$
\gamma_{right} :
g^* \circ a \to g^* \circ a \circ h_* \circ h^*
\xleftarrow{\xi_{right}} g^* \circ g_* \circ a' \circ h^* \to a' \circ h^*
$$
où $\xi_{right} : g_* \circ a' \to a \circ h_*$
est un isomorphisme (et peut donc être inversé)
et est la flèche « duale » du morphisme de changement de base
$h^* \circ f_* \to f'_* \circ g^*$. Les flèches extérieures proviennent
des morphismes canoniques $1 \to h_* \circ h^*$ et $g^* \circ g_* \to 1$.
De même, pour le carré de gauche, nous avons
$$
\gamma_{left} :
(g')^* \circ a' \to (g')^* \circ a' \circ (h')_* \circ (h')^*
\xleftarrow{\xi_{left}}
(g')^* \circ (g')_* \circ a'' \circ (h')^* \to a'' \circ (h')^*
$$
Pour le rectangle extérieur, nous obtenons
$$
\gamma_{rect} :
k^* \circ a \to
k^* \circ a \circ m_* \circ m^* \xleftarrow{\xi_{rect}}
k^* \circ k_* \circ a'' \circ m^* \to
a'' \circ m^*
$$
où $k = g \circ g'$ et $m = h \circ h'$.
Nous avons $k^* = (g')^* \circ g^*$ et $m^* = (h')^* \circ h^*$.
L'assertion du lemme est que $\gamma_{rect}$
est égal au composé
$$
k^* \circ a =
(g')^* \circ g^* \circ a \xrightarrow{\gamma_{right}}
(g')^* \circ a' \circ h^* \xrightarrow{\gamma_{left}}
a'' \circ (h')^* \circ h^* = a'' \circ m^*
$$
Pour le voir, considérons le diagramme suivant :
$$
\xymatrix{
& (g')^* \circ g^* \circ a \ar[d] \ar[ddl] \\
& (g')^* \circ g^* \circ a \circ h_* \circ h^* \ar[ld] \\
(g')^* \circ g^* \circ a \circ h_* \circ (h')_* \circ (h')^* \circ h^* &
(g')^* \circ g^* \circ g_* \circ a' \circ h^*
\ar[u]_{\xi_{right}} \ar[d] \ar[ld] \\
(g')^* \circ g^* \circ g_* \circ a' \circ (h')_* \circ (h')^* \circ h^*
\ar[u]_{\xi_{right}} \ar[dr] &
(g')^* \circ a' \circ h^* \ar[d] \\
(g')^* \circ g^* \circ g_* \circ (g')_* \circ a'' \circ (h')^* \circ h^*
\ar[u]_{\xi_{left}} \ar[ddr] \ar[dr] &
(g')^* \circ a' \circ (h')_* \circ (h')^* \circ h^* \\
& (g')^*\circ (g')_* \circ a'' \circ (h')^* \circ h^*
\ar[u]_{\xi_{left}} \ar[d] \\
& a'' \circ (h')^* \circ h^*
}
$$
En descendant du côté droit, nous obtenons le composé, et en descendant
du côté gauche, nous obtenons $\gamma_{rect}$.
Tous les quadrilatères du côté droit de ce diagramme sont commutatifs
d'après Catégories, Lemme \ref{categories-lemma-properties-2-cat-cats},
ou, plus simplement, d'après la discussion qui précède
Catégories, Définition \ref{categories-definition-horizontal-composition}.
Nous voyons donc qu'il suffit de montrer que
$$
g_* \circ (g')_* \circ a'' \xrightarrow{\xi_{left}}
g_* \circ a' \circ (h')_* \xrightarrow{\xi_{right}}
a \circ h_* \circ (h')_*
$$
est égal à $\xi_{rect}$. C'est l'assertion duale de
Cohomologie, Remarque \ref{cohomology-remark-compose-base-change-horizontal},
ce qui achève la démonstration.
\end{proof}

\begin{remark}
\label{remark-going-around}
Soit $S$ un schéma. Considérons un diagramme commutatif
$$
\xymatrix{
X'' \ar[r]_{k'} \ar[d]_{f''} & X' \ar[r]_k \ar[d]_{f'} & X \ar[d]^f \\
Y'' \ar[r]^{l'} \ar[d]_{g''} & Y' \ar[r]^l \ar[d]_{g'} & Y \ar[d]^g \\
Z'' \ar[r]^{m'} & Z' \ar[r]^m & Z
}
$$
d'espaces algébriques quasi-compacts et quasi-séparés sur $S$, dans lequel
tous les carrés sont cartésiens et où
$(f, l)$, $(g, m)$, $(f', l')$, $(g', m')$ sont des
paires de morphismes indépendants au sens de Tor.
Soient $a$, $a'$, $a''$, $b$, $b'$ et $b''$ les
adjoints à droite du Lemme \ref{lemma-twisted-inverse-image}
pour $f$, $f'$, $f''$, $g$, $g'$ et $g''$.
Étiquetons les carrés du diagramme $A$, $B$, $C$, $D$
comme suit :
$$
\begin{matrix}
A & B \\
C & D
\end{matrix}
$$
Les morphismes (\ref{equation-base-change-map})
des carrés sont alors les suivants (où nous utilisons $k^* = Lk^*$, etc.) :
$$
\begin{matrix}
\gamma_A : (k')^* \circ a' \to a'' \circ (l')^* &
\gamma_B : k^* \circ a \to a' \circ l^* \\
\gamma_C : (l')^* \circ b' \to b'' \circ (m')^* &
\gamma_D : l^* \circ b \to b' \circ m^*
\end{matrix}
$$
Pour les rectangles $2 \times 1$ et $1 \times 2$, nous avons quatre autres
morphismes de changement de base :
$$
\begin{matrix}
\gamma_{A + B} : (k \circ k')^* \circ a \to a'' \circ (l \circ l')^* \\
\gamma_{C + D} : (l \circ l')^* \circ b \to b'' \circ (m \circ m')^* \\
\gamma_{A + C} : (k')^* \circ (a' \circ b') \to (a'' \circ b'') \circ (m')^* \\
\gamma_{A + C} : k^* \circ (a \circ b) \to (a' \circ b') \circ m^*
\end{matrix}
$$
D'après le Lemme \ref{lemma-compose-base-change-maps-horizontal}, nous avons
$$
\gamma_{A + B} = \gamma_A \circ \gamma_B, \quad
\gamma_{C + D} = \gamma_C \circ \gamma_D
$$
et d'après le Lemme \ref{lemma-compose-base-change-maps}, nous avons
$$
\gamma_{A + C} = \gamma_C \circ \gamma_A, \quad
\gamma_{B + D} = \gamma_D \circ \gamma_B
$$
Il serait plus correct d'écrire ici
$\gamma_{A + B} = (\gamma_A \star \text{id}_{l^*}) \circ
(\text{id}_{(k')^*} \star \gamma_B)$ avec les notations de
Catégories, section \ref{categories-section-formal-cat-cat},
et de même pour les autres. Nous continuons toutefois à employer
l'abus de notation des démonstrations des
Lemmes \ref{lemma-compose-base-change-maps} et
\ref{lemma-compose-base-change-maps-horizontal},
qui consiste à omettre les produits $\star$ avec les identités, car on peut déterminer
ceux qu'il faut ajouter dès que la source et le but de la
transformation sont connus.
Cela étant dit, nous trouvons (a priori) deux transformations
$$
(k')^* \circ k^* \circ a \circ b
\longrightarrow
a'' \circ b'' \circ (m')^* \circ m^*
$$
à savoir
$$
\gamma_C \circ \gamma_A \circ \gamma_D \circ \gamma_B =
\gamma_{A + C} \circ \gamma_{B + D}
$$
et
$$
\gamma_C \circ \gamma_D \circ \gamma_A \circ \gamma_B =
\gamma_{C + D} \circ \gamma_{A + B}
$$
Le but de cette remarque est de signaler que ces transformations
sont égales. Pour le voir, il suffit en effet de montrer que
$$
\xymatrix{
(k')^* \circ a' \circ l^* \circ b \ar[r]_{\gamma_D} \ar[d]_{\gamma_A} &
(k')^* \circ a' \circ b' \circ m^* \ar[d]^{\gamma_A} \\
a'' \circ (l')^* \circ l^* \circ b \ar[r]^{\gamma_D} &
a'' \circ (l')^* \circ b' \circ m^*
}
$$
est commutatif. Cela résulte de
Catégories, Lemme \ref{categories-lemma-properties-2-cat-cats},
ou, plus simplement, de la discussion qui précède
Catégories, Définition \ref{categories-definition-horizontal-composition}.
\end{remark}







 
\section{Adjoint à droite de l'image directe et changement de base, II}
\label{section-base-change-II}

\noindent
Dans cette section, nous démontrons que le morphisme de changement de base de la
section \ref{section-base-change-map} est un isomorphisme
dans certains cas.

\begin{lemma}
\label{lemma-more-base-change}
Dans le diagramme (\ref{equation-base-change}), supposons en outre que
$g : Y' \to Y$ soit un morphisme de schémas affines et que $f : X \to Y$ soit propre.
Alors le morphisme de changement de base (\ref{equation-base-change-map}) induit un
isomorphisme
$$
L(g')^*a(K) \longrightarrow a'(Lg^*K)
$$
dans les cas suivants :
\begin{enumerate}
\item pour tout $K \in D_\QCoh(\mathcal{O}_X)$ si $f$
est plat et de présentation finie,
\item pour tout $K \in D_\QCoh(\mathcal{O}_X)$ si $f$
est parfait et $Y$ noethérien,
\item pour $K \in D_\QCoh^+(\mathcal{O}_X)$ si $g$ est de dimension de Tor finie
et $Y$ est noethérien.
\end{enumerate}
\end{lemma}

\begin{proof}
Écrivons $Y = \Spec(A)$ et $Y' = \Spec(A')$. Comme changement de base d'un morphisme
affine, le morphisme $g'$ est affine. Soit $M$ un générateur parfait
de $D_\QCoh(\mathcal{O}_X)$, voir Catégories dérivées des espaces, Théorème
\ref{spaces-perfect-theorem-bondal-van-den-Bergh}. Alors $L(g')^*M$ est un
générateur de $D_\QCoh(\mathcal{O}_{X'})$, voir
Catégories dérivées des espaces, Remarque
\ref{spaces-perfect-remark-pullback-generator}.
Il suffit donc de montrer que (\ref{equation-base-change-map})
induit un isomorphisme
\begin{equation}
\label{equation-iso}
R\Hom_{X'}(L(g')^*M, L(g')^*a(K))
\longrightarrow
R\Hom_{X'}(L(g')^*M, a'(Lg^*K))
\end{equation}
de complexes de Hom globaux, voir
Cohomologie sur les sites, section \ref{sites-cohomology-section-global-RHom},
car cela impliquera que le cône de $L(g')^*a(K) \to a'(Lg^*K)$
est nul.
La démonstration est structurée comme suit : nous montrerons d'abord que
ces complexes de Hom sont isomorphes et, dans sa dernière partie,
que l'isomorphisme est induit par (\ref{equation-iso}).

\medskip\noindent
Le membre de gauche. Comme $M$ est parfait, le morphisme canonique
$$
R\Hom_X(M, a(K)) \otimes^\mathbf{L}_A A'
\longrightarrow
R\Hom_{X'}(L(g')^*M, L(g')^*a(K))
$$
est un isomorphisme d'après Catégories dérivées des espaces, Lemme
\ref{spaces-perfect-lemma-affine-morphism-and-hom-out-of-perfect}.
En le combinant avec l'isomorphisme
$R\Hom_Y(Rf_*M, K) = R\Hom_X(M, a(K))$
du Lemme \ref{lemma-iso-global-hom},
nous obtenons que le membre de gauche est égal à
$R\Hom_Y(Rf_*M, K) \otimes^\mathbf{L}_A A'$.

\medskip\noindent
Le membre de droite. Nous utilisons d'abord ici l'isomorphisme
$$
R\Hom_{X'}(L(g')^*M, a'(Lg^*K)) = R\Hom_{Y'}(Rf'_*L(g')^*M, Lg^*K)
$$
du Lemme \ref{lemma-iso-global-hom}. Comme $f$ et $g$ sont
indépendants au sens de Tor, le morphisme de changement de base
$Lg^*Rf_*M \to Rf'_*L(g')^*M$ est un isomorphisme d'après
Catégories dérivées des espaces, Lemme
\ref{spaces-perfect-lemma-compare-base-change}.
Nous pouvons donc le réécrire sous la forme $R\Hom_{Y'}(Lg^*Rf_*M, Lg^*K)$.
Comme $Y$, $Y'$ sont affines et que $K$, $Rf_*M$ appartiennent à $D_\QCoh(\mathcal{O}_Y)$
(Catégories dérivées des espaces, Lemme
\ref{spaces-perfect-lemma-quasi-coherence-direct-image}),
nous avons un morphisme canonique
$$
\beta :
R\Hom_Y(Rf_*M, K) \otimes^\mathbf{L}_A A'
\longrightarrow
R\Hom_{Y'}(Lg^*Rf_*M, Lg^*K)
$$
dans $D(A')$. C'est la flèche de
Compléments d'algèbre, Équation (\ref{more-algebra-equation-base-change-RHom}),
où nous avons utilisé Catégories dérivées des schémas, Lemmes
\ref{perfect-lemma-affine-compare-bounded} et
\ref{perfect-lemma-quasi-coherence-internal-hom},
pour passer de la géométrie à l'algèbre et réciproquement.
\begin{enumerate}
\item Si $f$ est plat et de présentation finie, le complexe $Rf_*M$
est parfait sur $Y$ d'après Catégories dérivées des espaces, Lemme
\ref{spaces-perfect-lemma-flat-proper-perfect-direct-image-general},
et $\beta$ est un isomorphisme d'après
Compléments d'algèbre, Lemme \ref{more-algebra-lemma-base-change-RHom}, partie (1).
\item Si $f$ est parfait et $Y$ noethérien, le complexe $Rf_*M$
est parfait sur $Y$ d'après Compléments sur les morphismes d'espaces, Lemme
\ref{spaces-more-morphisms-lemma-perfect-proper-perfect-direct-image},
et $\beta$ est un isomorphisme comme précédemment.
\item Si $g$ est de dimension de Tor finie et $Y$ est noethérien,
le complexe $Rf_*M$ est pseudo-cohérent sur $Y$
(Catégories dérivées des espaces, Lemmes
\ref{spaces-perfect-lemma-direct-image-coherent} et
\ref{spaces-perfect-lemma-identify-pseudo-coherent-noetherian}),
et $\beta$ est un isomorphisme d'après
Compléments d'algèbre, Lemme \ref{more-algebra-lemma-base-change-RHom}, partie (4).
\end{enumerate}
Nous concluons que nous obtenons le même résultat que dans le paragraphe précédent.

\medskip\noindent
Dans la suite de la démonstration, nous montrons que les identifications des
membres de gauche et de droite de (\ref{equation-iso})
données dans les deuxième et troisième paragraphes sont effectivement données par
(\ref{equation-iso}). Pour que nos formules restent maniables,
nous noterons $(-, -)_X = R\Hom_X(-, -)$, écrirons $- \otimes A'$
à la place de $- \otimes_A^\mathbf{L} A'$, et abrégerons
$g^* = Lg^*$ et $f_* = Rf_*$. Considérons le diagramme
commutatif suivant :
$$
\xymatrix{
((g')^*M, (g')^*a(K))_{X'} \ar[d] &
(M, a(K))_X \otimes A' \ar[l]^-\alpha \ar[d] &
(f_*M, K)_Y \otimes A' \ar@{=}[l] \ar[d] \\
((g')^*M, (g')^*a(g_*g^*K))_{X'} &
(M, a(g_*g^*K))_X \otimes A' \ar[l]^-\alpha &
(f_*M, g_*g^*K)_Y \otimes A' \ar@{=}[l] \ar@/_4pc/[dd]_{\mu'} \\
((g')^*M, (g')^*g'_*a'(g^*K))_{X'} \ar[u] \ar[d] &
(M, g'_*a'(g^*K))_X \otimes A' \ar[u] \ar[l]^-\alpha \ar[ld]^\mu &
(f_*M, K) \otimes A' \ar[d]^\beta \\
((g')^*M, a'(g^*K))_{X'} &
(f'_*(g')^*M, g^*K)_{Y'} \ar@{=}[l] \ar[r] &
(g^*f_*M, g^*K)_{Y'}
}
$$
Les flèches étiquetées $\alpha$ sont les morphismes de
Catégories dérivées des espaces, Lemme
\ref{spaces-perfect-lemma-affine-morphism-and-hom-out-of-perfect},
pour le diagramme de sommets $X', X, Y', Y$.
La partie supérieure du diagramme est commutative, car les flèches horizontales sont
fonctorielles en leurs arguments.
Les flèches verticales du milieu proviennent de la transformation inversible
$g'_* \circ a' \to a \circ g_*$ du Lemme \ref{lemma-flat-precompose-pus},
et le carré médian est donc commutatif.
Descendre le long du côté gauche donne (\ref{equation-iso}).
Les flèches horizontales supérieures donnent les identifications utilisées dans le
deuxième paragraphe de la démonstration.
Les flèches horizontales inférieures, y compris $\beta$, donnent les identifications
utilisées dans le troisième paragraphe de la démonstration. Étant donnés $E \in D(A)$,
$E' \in D(A')$ et $c : E \to E'$ dans $D(A)$, nous noterons
$\mu_c : E \otimes A' \to E'$ le morphisme induit par $c$
et par l'adjonction entre restriction et changement de base ;
si $c$ est clair, nous écrivons $\mu = \mu_c$, c'est-à-dire que nous
omettons $c$ de la notation. Le morphisme $\mu$ du diagramme est de cette
forme, avec $c$ donné par l'identification
$(M, g'_*a(g^*K))_X = ((g')^*M, a'(g^*K))_{X'}$
; le triangle contenant $\mu$ est commutatif d'après
Catégories dérivées des espaces, Remarque
\ref{spaces-perfect-remark-multiplication-map}.

\medskip\noindent
Observons que
$$
\xymatrix{
(M, a(g_*g^*K))_X &
(f_*M, g_* g^*K)_Y \ar@{=}[l] &
(g^*f_*M, g^*K)_{Y'} \ar@{=}[l] \\
(M, g'_* a'(g^*K))_X \ar[u] &
((g')^*M, a'(g^*K))_{X'} \ar@{=}[l] &
(f'_*(g')^*M, g^*K)_{Y'} \ar@{=}[l] \ar[u]
}
$$
est commutatif par la définition même de la transformation
$g'_* \circ a' \to a \circ g_*$. Si $\mu'$ désigne, comme ci-dessus, le morphisme
correspondant à l'identification
$(f_*M, g_*g^*K)_X = (g^*f_*M, g^*K)_{Y'}$, alors
l'hexagone est lui aussi commutatif. Il suffit donc de montrer que
$\beta$ est égal au composé de
$(f_*M, K)_Y \otimes A' \to (f_*M, g_*g^*K)_X \otimes A'$
et de $\mu'$. Pour cela, il suffit de démontrer que les deux morphismes induits
$(f_*M, K)_Y \to (g^*f_*M, g^*K)_{Y'}$ sont égaux.
Autrement dit, il suffit de montrer que le diagramme
$$
\xymatrix{
R\Hom_A(E, K) \ar[rr]_{\text{induit par }\beta} \ar[rd] & &
R\Hom_{A'}(E \otimes_A^\mathbf{L} A', K \otimes_A^\mathbf{L} A') \\
& R\Hom_A(E, K \otimes_A^\mathbf{L} A') \ar[ru]
}
$$
est commutatif pour tous $E, K \in D(A)$. Puisque c'est ainsi que $\beta$ est construit dans
Compléments d'algèbre, section \ref{more-algebra-section-base-change-RHom},
la démonstration est achevée.
\end{proof}





 
\section{Adjoint à droite de l'image directe et morphismes trace}
\label{section-trace}

\noindent
Soit $S$ un schéma.
Soit $f : X \to Y$ un morphisme d'espaces algébriques quasi-compacts et quasi-séparés
sur $S$.
Soit $a : D_\QCoh(\mathcal{O}_Y) \to D_\QCoh(\mathcal{O}_X)$
l'adjoint à droite du Lemme \ref{lemma-twisted-inverse-image}. D'après
Catégories, section \ref{categories-section-adjoint}, nous obtenons une
transformation de foncteurs
$$
\text{Tr}_f : Rf_* \circ a \longrightarrow \text{id}
$$
Le morphisme correspondant $\text{Tr}_{f, K} : Rf_*a(K) \longrightarrow K$
pour $K \in D_\QCoh(\mathcal{O}_Y)$ est parfois appelé le {\it morphisme trace}.
C'est le morphisme ayant la propriété que la bijection
$$
\Hom_X(L, a(K)) \longrightarrow \Hom_Y(Rf_*L, K)
$$
pour $L \in D_\QCoh(\mathcal{O}_X)$ qui caractérise l'adjoint à droite
est donnée par
$$
\varphi \longmapsto \text{Tr}_{f, K} \circ Rf_*\varphi
$$
Le morphisme canonique (\ref{equation-sheafy-trace})
$$
Rf_*R\SheafHom_{\mathcal{O}_X}(L, a(K))
\longrightarrow
R\SheafHom_{\mathcal{O}_Y}(Rf_*L, K)
$$
est obtenu par composition avec $\text{Tr}_{f, K}$.
Tout morphisme trace que nous considérerons dans cette section sera un
cas particulier de celui-ci. Avant d'en étudier quelques cas particuliers,
nous montrons que la formation du morphisme trace commute au changement de base.

\begin{lemma}[Morphisme trace et changement de base]
\label{lemma-trace-map-and-base-change}
Supposons donné un diagramme (\ref{equation-base-change}).
Alors les morphismes
$1 \star \text{Tr}_f : Lg^* \circ Rf_* \circ a \to Lg^*$ et
$\text{Tr}_{f'} \star 1 : Rf'_* \circ a' \circ Lg^* \to Lg^*$
coïncident via les morphismes de changement de base
$\beta : Lg^* \circ Rf_* \to Rf'_* \circ L(g')^*$
(Cohomologie sur les sites, Remarque \ref{sites-cohomology-remark-base-change})
et $\alpha : L(g')^* \circ a \to a' \circ Lg^*$
(\ref{equation-base-change-map}).
Plus précisément, le diagramme
$$
\xymatrix{
Lg^* \circ Rf_* \circ a
\ar[d]_{\beta \star 1} \ar[r]_-{1 \star \text{Tr}_f} &
Lg^* \\
Rf'_* \circ L(g')^* \circ a \ar[r]^{1 \star \alpha} &
Rf'_* \circ a' \circ Lg^* \ar[u]_{\text{Tr}_{f'} \star 1}
}
$$
de transformations de foncteurs est commutatif.
\end{lemma}

\begin{proof}
Dans cette démonstration, nous écrivons $f_*$ pour $Rf_*$ et $g^*$ pour $Lg^*$, et nous
omettons les produits $\star$ avec les identités, puisqu'on peut déterminer ceux qu'il faut
ajouter dès que la source et le but de la transformation sont connus.
Rappelons que $\beta : g^* \circ f_* \to f'_* \circ (g')^*$ est un isomorphisme
et que $\alpha$ est défini à l'aide de
l'isomorphisme $\beta^\vee : g'_* \circ a' \to a \circ g_*$,
qui est l'adjoint de $\beta$ ; voir le Lemme \ref{lemma-flat-precompose-pus}
et sa démonstration. Remarquons d'abord que la flèche horizontale supérieure
du diagramme du lemme est égale au composé
$$
g^* \circ f_* \circ a \to
g^* \circ f_* \circ a \circ g_* \circ g^* \to
g^* \circ g_* \circ g^* \to g^*
$$
où la première flèche est l'unité de $(g^*, g_*)$, la deuxième flèche
est $\text{Tr}_f$, et la troisième est la co-unité de $(g^*, g_*)$.
C'est une conséquence immédiate du fait que le composé
$g^* \to g^* \circ g_* \circ g^* \to g^*$ de l'unité et de la co-unité est l'identité.
Considérons le diagramme
$$
\xymatrix{
& g^* \circ f_* \circ a \ar[ld]_\beta \ar[d] \ar[r]_{\text{Tr}_f} & g^* \\
f'_* \circ (g')^* \circ a \ar[dr] &
g^* \circ f_* \circ a \circ g_* \circ g^* \ar[d]_\beta \ar[ru] &
g^* \circ f_* \circ g'_* \circ a' \circ g^* \ar[l]_{\beta^\vee} \ar[d]_\beta &
f'_* \circ a' \circ g^* \ar[lu]_{\text{Tr}_{f'}} \\
& f'_* \circ (g')^* \circ a \circ g_* \circ g^* &
f'_* \circ (g')^* \circ g'_* \circ a' \circ g^* \ar[ru] \ar[l]_{\beta^\vee}
}
$$
Dans ce diagramme, les deux carrés sont commutatifs d'après
Catégories, Lemme \ref{categories-lemma-properties-2-cat-cats},
ou, plus simplement, d'après la discussion qui précède
Catégories, Définition \ref{categories-definition-horizontal-composition}.
Le triangle est commutatif d'après la discussion ci-dessus. D'après
Catégories, Lemme
\ref{categories-lemma-transformation-between-functors-and-adjoints},
le carré
$$
\xymatrix{
g^* \circ f_* \circ g'_* \circ a' \ar[d]_{\beta^\vee} \ar[r]_-\beta &
f'_* \circ (g')^* \circ g'_* \circ a' \ar[d] \\
g^* \circ f_* \circ a \circ g_* \ar[r] &
\text{id}
}
$$
est commutatif, ce qui implique que le pentagone du grand diagramme est commutatif.
Comme $\beta$ et $\beta^\vee$ sont des isomorphismes, et comme le parcours
du bord extérieur du grand diagramme est égal à
$\text{Tr}_f \circ \alpha \circ \beta$ par définition, cela démontre le lemme.
\end{proof}

\noindent
Soit $S$ un schéma.
Soit $f : X \to Y$ un morphisme d'espaces algébriques quasi-compacts et quasi-séparés
sur $S$.
Soit $a : D_\QCoh(\mathcal{O}_Y) \to D_\QCoh(\mathcal{O}_X)$
l'adjoint à droite de $Rf_*$ comme dans le
Lemme \ref{lemma-twisted-inverse-image}. D'après
Catégories, section \ref{categories-section-adjoint}, nous obtenons une
transformation de foncteurs
$$
\eta_f : \text{id} \to  a \circ Rf_*
$$
appelée l'unité de l'adjonction.

\begin{lemma}
\label{lemma-unit-and-base-change}
Supposons donné un diagramme (\ref{equation-base-change}). Alors les morphismes
$1 \star \eta_f : L(g')^* \to L(g')^* \circ a \circ Rf_*$ et
$\eta_{f'} \star 1 : L(g')^* \to a' \circ Rf'_* \circ L(g')^*$
coïncident via les morphismes de changement de base
$\beta : Lg^* \circ Rf_* \to Rf'_* \circ L(g')^*$
(Cohomologie sur les sites, Remarque \ref{sites-cohomology-remark-base-change})
et $\alpha : L(g')^* \circ a \to a' \circ Lg^*$
(\ref{equation-base-change-map}).
Plus précisément, le diagramme
$$
\xymatrix{
L(g')^* \ar[r]_-{1 \star \eta_f} \ar[d]_{\eta_{f'} \star 1} &
L(g')^* \circ a \circ Rf_* \ar[d]^\alpha \\
a' \circ Rf'_* \circ L(g')^* &
a' \circ Lg^* \circ Rf_* \ar[l]_-\beta
}
$$
de transformations de foncteurs est commutatif.
\end{lemma}

\begin{proof}
Cette démonstration est duale de celle du Lemme \ref{lemma-trace-map-and-base-change}.
Dans cette démonstration, nous écrivons $f_*$ pour $Rf_*$ et $g^*$ pour $Lg^*$, et nous
omettons les produits $\star$ avec les identités, puisqu'on peut déterminer ceux qu'il faut
ajouter dès que la source et le but de la transformation sont connus.
Rappelons que $\beta : g^* \circ f_* \to f'_* \circ (g')^*$ est un isomorphisme
et que $\alpha$ est défini à l'aide de
l'isomorphisme $\beta^\vee : g'_* \circ a' \to a \circ g_*$,
qui est l'adjoint de $\beta$ ; voir le Lemme \ref{lemma-flat-precompose-pus}
et sa démonstration. Remarquons d'abord que la flèche verticale gauche
du diagramme du lemme est égale au composé
$$
(g')^* \to (g')^* \circ g'_* \circ (g')^* \to
(g')^* \circ g'_* \circ a' \circ f'_* \circ (g')^* \to
a' \circ f'_* \circ (g')^*
$$
où la première flèche est l'unité de $((g')^*, g'_*)$, la deuxième flèche
est $\eta_{f'}$, et la troisième est la co-unité de $((g')^*, g'_*)$.
C'est une conséquence immédiate du fait que le composé
$(g')^* \to (g')^* \circ (g')_* \circ (g')^* \to (g')^*$
de l'unité et de la co-unité est l'identité. Considérons le diagramme
$$
\xymatrix{
& (g')^* \circ a \circ f_* \ar[r] &
(g')^* \circ a \circ g_* \circ g^* \circ f_*
\ar[ld]_\beta \\
(g')^* \ar[ru]^{\eta_f} \ar[dd]_{\eta_{f'}} \ar[rd] &
(g')^* \circ a \circ g_* \circ f'_* \circ (g')^* &
(g')^* \circ g'_* \circ a' \circ g^* \circ f_*
\ar[u]_{\beta^\vee} \ar[ld]_\beta \ar[d] \\
& (g')^* \circ g'_* \circ a' \circ f'_* \circ (g')^*
\ar[ld] \ar[u]_{\beta^\vee} &
a' \circ g^* \circ f_* \ar[lld]^\beta \\
a' \circ f'_* \circ (g')^*
}
$$
Dans ce diagramme, les deux carrés sont commutatifs d'après
Catégories, Lemme \ref{categories-lemma-properties-2-cat-cats},
ou, plus simplement, d'après la discussion qui précède
Catégories, Définition \ref{categories-definition-horizontal-composition}.
Le triangle est commutatif d'après la discussion ci-dessus. D'après le dual de
Catégories, Lemme
\ref{categories-lemma-transformation-between-functors-and-adjoints},
le carré
$$
\xymatrix{
\text{id} \ar[r] \ar[d] &
g'_* \circ a' \circ g^* \circ f_* \ar[d]^\beta \\
g'_* \circ a' \circ g^* \circ f_* \ar[r]^{\beta^\vee} &
a \circ g_* \circ f'_* \circ (g')^*
}
$$
est commutatif, ce qui implique que le pentagone du grand diagramme est commutatif.
Comme $\beta$ et $\beta^\vee$ sont des isomorphismes, et comme le parcours
du bord extérieur du grand diagramme est égal à
$\beta \circ \alpha \circ \eta_f$ par définition, cela démontre le lemme.
\end{proof}





 
\section{Adjoint à droite de l'image directe et image réciproque}
\label{section-compare-with-pullback}

\noindent
Soit $S$ un schéma.
Soit $f : X \to Y$ un morphisme d'espaces algébriques quasi-compacts et quasi-séparés
sur $S$.
Soit $a$ l'adjoint à droite de l'image directe, comme dans le
Lemme \ref{lemma-twisted-inverse-image}. Pour $K, L \in D_\QCoh(\mathcal{O}_Y)$,
il existe un morphisme canonique
$$
Lf^*K \otimes^\mathbf{L}_{\mathcal{O}_X} a(L)
\longrightarrow
a(K \otimes_{\mathcal{O}_Y}^\mathbf{L} L)
$$
Ce morphisme est en effet adjoint d'un morphisme
$$
Rf_*(Lf^*K \otimes^\mathbf{L}_{\mathcal{O}_X} a(L)) =
K \otimes^\mathbf{L}_{\mathcal{O}_Y} Rf_*(a(L))
\longrightarrow
K \otimes^\mathbf{L}_{\mathcal{O}_Y} L
$$
(l'égalité résulte de Catégories dérivées des espaces, Lemme
\ref{spaces-perfect-lemma-cohomology-base-change}),
pour lequel nous utilisons le morphisme trace $Rf_*a(L) \to L$.
Lorsque $L = \mathcal{O}_Y$, nous obtenons un morphisme
\begin{equation}
\label{equation-compare-with-pullback}
Lf^*K \otimes^\mathbf{L}_{\mathcal{O}_X} a(\mathcal{O}_Y) \longrightarrow a(K)
\end{equation}
fonctoriel en $K$ et compatible aux triangles distingués.

\begin{lemma}
\label{lemma-compare-with-pullback-perfect}
Soit $S$ un schéma.
Soit $f : X \to Y$ un morphisme d'espaces algébriques quasi-compacts et quasi-séparés
sur $S$. Le morphisme
$Lf^*K \otimes^\mathbf{L}_{\mathcal{O}_X} a(L) \to
a(K \otimes_{\mathcal{O}_Y}^\mathbf{L} L)$
défini ci-dessus pour $K, L \in D_\QCoh(\mathcal{O}_Y)$
est un isomorphisme si $K$ est parfait. En particulier,
(\ref{equation-compare-with-pullback}) est un isomorphisme si $K$ est parfait.
\end{lemma}

\begin{proof}
Soit $K^\vee$ le « dual » de $K$, voir
Cohomologie sur les sites, Lemme \ref{sites-cohomology-lemma-dual-perfect-complex}.
Pour $M \in D_\QCoh(\mathcal{O}_X)$, nous avons
\begin{align*}
\Hom_{D(\mathcal{O}_Y)}(Rf_*M, K \otimes^\mathbf{L}_{\mathcal{O}_Y} L)
& =
\Hom_{D(\mathcal{O}_Y)}(
Rf_*M \otimes^\mathbf{L}_{\mathcal{O}_Y} K^\vee, L) \\
& =
\Hom_{D(\mathcal{O}_X)}(
M \otimes^\mathbf{L}_{\mathcal{O}_X} Lf^*K^\vee, a(L)) \\
& =
\Hom_{D(\mathcal{O}_X)}(M,
Lf^*K \otimes^\mathbf{L}_{\mathcal{O}_X} a(L))
\end{align*}
La deuxième égalité résulte de la définition de $a$ et de la formule de projection
(Cohomologie sur les sites, Lemme
\ref{sites-cohomology-lemma-projection-formula}),
ou du résultat plus général Catégories dérivées des espaces, Lemme
\ref{spaces-perfect-lemma-cohomology-base-change}.
Le résultat découle donc du lemme de Yoneda.
\end{proof}

\begin{lemma}
\label{lemma-restriction-compare-with-pullback}
Supposons donné un diagramme (\ref{equation-base-change}).
Soit $K \in D_\QCoh(\mathcal{O}_Y)$. Le diagramme
$$
\xymatrix{
L(g')^*(Lf^*K \otimes^\mathbf{L}_{\mathcal{O}_X} a(\mathcal{O}_Y))
\ar[r] \ar[d] & L(g')^*a(K) \ar[d] \\
L(f')^*Lg^*K \otimes_{\mathcal{O}_{X'}}^\mathbf{L} a'(\mathcal{O}_{Y'})
\ar[r] & a'(Lg^*K)
}
$$
est commutatif, où les flèches horizontales sont les morphismes
(\ref{equation-compare-with-pullback}) pour $K$ et $Lg^*K$,
et où les flèches verticales sont construites à l'aide de
Cohomologie sur les sites, Remarque \ref{sites-cohomology-remark-base-change}, et de
(\ref{equation-base-change-map}).
\end{lemma}

\begin{proof}
Dans cette démonstration, nous écrirons $f_*$ pour $Rf_*$ et $f^*$ pour $Lf^*$, etc.,
et nous écrirons $\otimes$ pour $\otimes^\mathbf{L}_{\mathcal{O}_X}$, etc.
Écrivons (\ref{equation-compare-with-pullback}) comme le composé
\begin{align*}
f^*K \otimes a(\mathcal{O}_Y)
& \to
a(f_*(f^*K \otimes a(\mathcal{O}_Y))) \\
& \leftarrow
a(K \otimes f_*a(\mathcal{O}_K)) \\
& \to
a(K \otimes \mathcal{O}_Y) \\
& \to
a(K)
\end{align*}
Ici, la première flèche est l'unité $\eta_f$, la deuxième est l'image par $a$
de Cohomologie sur les sites, Équation
(\ref{sites-cohomology-equation-projection-formula-map}), qui est un
isomorphisme d'après Catégories dérivées des espaces, Lemme
\ref{spaces-perfect-lemma-cohomology-base-change} ; la troisième flèche est
l'image par $a$ de $\text{id}_K \otimes \text{Tr}_f$, et la quatrième
est l'image par $a$ de l'isomorphisme $K \otimes \mathcal{O}_Y = K$.
La démonstration du lemme consiste à montrer que chacun de ces
morphismes donne un carré commutatif comme dans l'énoncé du lemme.
Pour $\eta_f$ et $\text{Tr}_f$, cela résulte des
Lemmes \ref{lemma-unit-and-base-change} et
\ref{lemma-trace-map-and-base-change}.
Pour la flèche qui utilise Cohomologie sur les sites, Équation
(\ref{sites-cohomology-equation-projection-formula-map}),
cela résulte de Cohomologie sur les sites, Remarque
\ref{sites-cohomology-remark-compatible-with-diagram}.
Pour le morphisme de multiplication, c'est clair. Cela achève la démonstration.
\end{proof}









 
\section{Adjoint à droite de l'image directe pour les morphismes propres et plats}
\label{section-proper-flat}

\noindent
Pour les morphismes propres, plats et de présentation finie entre espaces algébriques
quasi-compacts et quasi-séparés, l'adjoint à droite de l'image directe
possède quelques propriétés remarquables.

\begin{lemma}
\label{lemma-proper-flat}
Soit $S$ un schéma.
Soit $Y$ un espace algébrique quasi-compact et quasi-séparé sur $S$.
Soit $f : X \to Y$ un morphisme d'espaces algébriques propre, plat et
de présentation finie.
Soit $a$ l'adjoint à droite de
$Rf_* : D_\QCoh(\mathcal{O}_X) \to D_\QCoh(\mathcal{O}_Y)$ du
Lemme \ref{lemma-twisted-inverse-image}. Alors $a$ commute aux sommes directes.
\end{lemma}

\begin{proof}
Soit $P$ un objet parfait de $D(\mathcal{O}_X)$. D'après
Catégories dérivées des espaces, Lemme
\ref{spaces-perfect-lemma-flat-proper-perfect-direct-image-general},
le complexe $Rf_*P$ est parfait sur $Y$.
Soit $K_i$ une famille d'objets de $D_\QCoh(\mathcal{O}_Y)$.
Alors
\begin{align*}
\Hom_{D(\mathcal{O}_X)}(P, a(\bigoplus K_i))
& =
\Hom_{D(\mathcal{O}_Y)}(Rf_*P, \bigoplus K_i) \\
& =
\bigoplus \Hom_{D(\mathcal{O}_Y)}(Rf_*P, K_i) \\
& =
\bigoplus \Hom_{D(\mathcal{O}_X)}(P, a(K_i))
\end{align*}
car un objet parfait est compact (Catégories dérivées des espaces,
Proposition \ref{spaces-perfect-proposition-compact-is-perfect}).
Comme $D_\QCoh(\mathcal{O}_X)$ possède un générateur parfait
(Catégories dérivées des espaces, Théorème
\ref{spaces-perfect-theorem-bondal-van-den-Bergh}),
nous concluons que le morphisme $\bigoplus a(K_i) \to a(\bigoplus K_i)$
est un isomorphisme, c'est-à-dire que $a$ commute aux sommes directes.
\end{proof}

\begin{lemma}
\label{lemma-compare-with-pullback-flat-proper}
Soit $S$ un schéma.
Soit $Y$ un espace algébrique quasi-compact et quasi-séparé sur $S$.
Soit $f : X \to Y$ un morphisme d'espaces algébriques propre, plat et
de présentation finie.
Le morphisme (\ref{equation-compare-with-pullback}) est un isomorphisme
pour tout objet $K$ de $D_\QCoh(\mathcal{O}_Y)$.
\end{lemma}

\begin{proof}
D'après le Lemme \ref{lemma-proper-flat}, nous savons que $a$ commute
aux sommes directes. La collection des objets de
$D_\QCoh(\mathcal{O}_Y)$ pour lesquels (\ref{equation-compare-with-pullback})
est un isomorphisme est donc une sous-catégorie strictement pleine, saturée et triangulée
de $D_\QCoh(\mathcal{O}_Y)$, qui est en outre
stable par sommes directes. Comme $D_\QCoh(\mathcal{O}_Y)$
est une catégorie de modules (Catégories dérivées des espaces, Théorème
\ref{spaces-perfect-theorem-DQCoh-is-Ddga}) engendrée par un unique
objet parfait (Catégories dérivées des espaces, Théorème
\ref{spaces-perfect-theorem-bondal-van-den-Bergh}),
nous pouvons raisonner comme dans
Compléments d'algèbre, Remarque \ref{more-algebra-remark-P-resolution},
pour voir qu'il suffit de démontrer que (\ref{equation-compare-with-pullback})
est un isomorphisme pour un seul objet parfait.
Or le résultat vaut pour les objets parfaits, voir
Lemme \ref{lemma-compare-with-pullback-perfect}.
\end{proof}

\begin{lemma}
\label{lemma-properties-relative-dualizing}
Soit $Y$ un schéma affine. Soit $f : X \to Y$ un morphisme
d'espaces algébriques propre, plat et de présentation finie.
Soit $a$ l'adjoint à droite de
$Rf_* : D_\QCoh(\mathcal{O}_X) \to D_\QCoh(\mathcal{O}_Y)$ du
Lemme \ref{lemma-twisted-inverse-image}.
Alors
\begin{enumerate}
\item $a(\mathcal{O}_Y)$ est un objet $Y$-parfait de $D(\mathcal{O}_X)$,
\item les faisceaux de cohomologie de $Rf_*a(\mathcal{O}_Y)$
sont nuls en degrés positifs,
\item $\mathcal{O}_X \to
R\SheafHom_{\mathcal{O}_X}(a(\mathcal{O}_Y), a(\mathcal{O}_Y))$
est un isomorphisme.
\end{enumerate}
\end{lemma}

\begin{proof}
Pour un objet parfait $E$ de $D(\mathcal{O}_X)$, nous avons
\begin{align*}
Rf_*(E \otimes_{\mathcal{O}_X}^\mathbf{L} \omega_{X/Y}^\bullet)
& =
Rf_*R\SheafHom_{\mathcal{O}_X}(E^\vee, \omega_{X/Y}^\bullet) \\
& =
R\SheafHom_{\mathcal{O}_Y}(Rf_*E^\vee, \mathcal{O}_Y) \\
& =
(Rf_*E^\vee)^\vee
\end{align*}
Pour la première égalité, voir
Cohomologie sur les sites, Lemme \ref{sites-cohomology-lemma-dual-perfect-complex}.
Pour la deuxième égalité, voir le Lemme \ref{lemma-iso-on-RSheafHom},
la Remarque \ref{remark-iso-on-RSheafHom}, et
Catégories dérivées des espaces, Lemme
\ref{spaces-perfect-lemma-flat-proper-perfect-direct-image-general}.
La troisième égalité est la définition du dual. En particulier,
ces références montrent aussi que le résultat est un objet parfait
de $D(\mathcal{O}_Y)$. Nous concluons que $\omega_{X/Y}^\bullet$
est $Y$-parfait d'après Compléments sur les morphismes d'espaces, Lemme
\ref{spaces-more-morphisms-lemma-characterize-relatively-perfect}.
Cela démontre (1).

\medskip\noindent
Soit $M$ un objet de $D_\QCoh(\mathcal{O}_Y)$. Alors
\begin{align*}
\Hom_Y(M, Rf_*a(\mathcal{O}_Y)) & =
\Hom_X(Lf^*M, a(\mathcal{O}_Y)) \\
& =
\Hom_Y(Rf_*Lf^*M, \mathcal{O}_Y) \\
& =
\Hom_Y(M \otimes_{\mathcal{O}_Y}^\mathbf{L} Rf_*\mathcal{O}_Y, \mathcal{O}_Y)
\end{align*}
La première égalité résulte de Cohomologie sur les sites, Lemme
\ref{sites-cohomology-lemma-adjoint}.
La deuxième égalité résulte de la construction de $a$.
La troisième égalité résulte de Catégories dérivées des espaces, Lemme
\ref{spaces-perfect-lemma-cohomology-base-change}.
Rappelons que $Rf_*\mathcal{O}_X$ est parfait, d'amplitude de Tor dans $[0, N]$
pour un certain $N$, voir
Catégories dérivées des espaces, Lemme
\ref{spaces-perfect-lemma-flat-proper-perfect-direct-image-general}.
Nous pouvons donc représenter $Rf_*\mathcal{O}_X$ par un complexe de
modules projectifs finis placé en degrés $[0, N]$
(en utilisant Compléments d'algèbre, Lemme \ref{more-algebra-lemma-perfect},
et le fait que $Y$ est affine).
Ainsi, si $M = \mathcal{O}_Y[-i]$ pour un certain $i > 0$, le dernier
groupe est nul. Comme $Y$ est affine, nous concluons que
$H^i(Rf_*a(\mathcal{O}_Y)) = 0$ pour $i > 0$.
Cela démontre (2).

\medskip\noindent
Soit $E$ un objet parfait de $D_\QCoh(\mathcal{O}_X)$. Alors
nous avons
\begin{align*}
\Hom_X(E, R\SheafHom_{\mathcal{O}_X}(a(\mathcal{O}_Y), a(\mathcal{O}_Y))
& =
\Hom_X(E \otimes_{\mathcal{O}_X}^\mathbf{L} a(\mathcal{O}_Y),
a(\mathcal{O}_Y)) \\
& =
\Hom_Y(Rf_*(E \otimes_{\mathcal{O}_X}^\mathbf{L} a(\mathcal{O}_Y)),
\mathcal{O}_Y) \\
& =
\Hom_Y(Rf_*(R\SheafHom_{\mathcal{O}_X}(E^\vee, a(\mathcal{O}_Y))),
\mathcal{O}_Y) \\
& =
\Hom_Y(R\SheafHom_{\mathcal{O}_Y}(Rf_*E^\vee, \mathcal{O}_Y),
\mathcal{O}_Y) \\
& =
R\Gamma(Y, Rf_*E^\vee) \\
& =
\Hom_X(E, \mathcal{O}_X)
\end{align*}
La première égalité résulte de Cohomologie sur les sites, Lemme
\ref{sites-cohomology-lemma-internal-hom}.
La deuxième égalité est la définition de $a$.
La troisième égalité provient de la construction du complexe parfait dual
$E^\vee$, voir Cohomologie sur les sites, Lemme
\ref{sites-cohomology-lemma-dual-perfect-complex}.
La quatrième égalité résulte de l'égalité
$Rf_*R\SheafHom_{\mathcal{O}_X}(E^\vee, \omega_{X/Y}^\bullet) =
R\SheafHom_{\mathcal{O}_Y}(Rf_*E^\vee, \mathcal{O}_Y)$
démontrée dans le premier paragraphe de la preuve.
La cinquième égalité résulte de la bidualité pour les complexes parfaits
(Cohomologie sur les sites, Lemme
\ref{sites-cohomology-lemma-dual-perfect-complex})
et du fait que $Rf_*E$ est parfait d'après
Catégories dérivées des espaces, Lemme
\ref{spaces-perfect-lemma-flat-proper-perfect-direct-image-general}.
La dernière égalité est la suite spectrale de Leray pour $f$.
Cette suite d'égalités montre essentiellement que (3)
vaut d'après le lemme de Yoneda. En effet, l'objet
$R\SheafHom(a(\mathcal{O}_Y), a(\mathcal{O}_Y))$
appartient à $D_\QCoh(\mathcal{O}_X)$ d'après Catégories dérivées des espaces, Lemme
\ref{spaces-perfect-lemma-quasi-coherence-internal-hom}.
En prenant $E = \mathcal{O}_X$ ci-dessus, nous obtenons un morphisme
$\alpha : \mathcal{O}_X \to
R\SheafHom_{\mathcal{O}_X}(a(\mathcal{O}_Y), a(\mathcal{O}_Y))$
correspondant à
$\text{id}_{\mathcal{O}_X} \in \Hom_X(\mathcal{O}_X, \mathcal{O}_X)$.
Comme tous les isomorphismes ci-dessus sont fonctoriels en $E$, nous
voyons que le cône de $\alpha$ est un objet $C$ de $D_\QCoh(\mathcal{O}_X)$
tel que $\Hom(E, C) = 0$ pour tout $E$ parfait.
Comme les objets parfaits engendrent
(Catégories dérivées des espaces, Théorème
\ref{spaces-perfect-theorem-bondal-van-den-Bergh}),
nous concluons que $\alpha$ est un isomorphisme.
\end{proof}









 
\section{Complexes dualisants relatifs pour les morphismes propres et plats}
\label{section-relative-dualizing-proper-flat}

\noindent
Motivés par Dualité pour les schémas, sections
\ref{duality-section-proper-flat} et
\ref{duality-section-relative-dualizing-complexes},
ainsi que par les résultats de la
section \ref{section-proper-flat},
nous posons la définition suivante.

\begin{definition}
\label{definition-relative-dualizing-proper-flat}
Soit $S$ un schéma. Soit $f : X \to Y$ un morphisme propre et plat
d'espaces algébriques sur $S$, de présentation finie.
Un {\it complexe dualisant relatif} pour $X/Y$ est un couple
$(\omega_{X/Y}^\bullet, \tau)$ formé d'un
objet $Y$-parfait $\omega_{X/Y}^\bullet$ de $D(\mathcal{O}_X)$
et d'un morphisme
$$
\tau : Rf_*\omega_{X/Y}^\bullet \longrightarrow \mathcal{O}_Y
$$
tel que, pour tout carré cartésien
$$
\xymatrix{
X' \ar[r]_{g'} \ar[d]_{f'} & X \ar[d]^f \\
Y' \ar[r]^g & Y
}
$$
où $Y'$ est un schéma affine, le couple
$(L(g')^*\omega_{X/Y}^\bullet, Lg^*\tau)$
soit isomorphe au couple
$(a'(\mathcal{O}_{Y'}), \text{Tr}_{f', \mathcal{O}_{Y'}})$
étudié dans les sections
\ref{section-twisted-inverse-image},
\ref{section-base-change-map},
\ref{section-base-change-II},
\ref{section-trace},
\ref{section-compare-with-pullback} et
\ref{section-proper-flat}.
\end{definition}

\noindent
Il convient de faire ici plusieurs remarques.
\begin{enumerate}
\item Dans la Définition \ref{definition-relative-dualizing-proper-flat},
on peut supprimer l'hypothèse que $\omega_{X/Y}^\bullet$ est $Y$-parfait.
En effet, en faisant parcourir à $Y'$ les membres d'un recouvrement \'etale de $Y$
par des affines, le Lemme \ref{lemma-properties-relative-dualizing} montre
que les restrictions de $\omega_{X/Y}^\bullet$ aux membres d'un
recouvrement \'etale de $X$ sont $Y$-parfaites, ce qui implique que
$\omega_{X/Y}^\bullet$ est $Y$-parfait ; voir
Compléments sur les morphismes d'espaces, section
\ref{spaces-more-morphisms-section-relatively-perfect}.
\item Considérons un complexe dualisant relatif
$(\omega_{X/Y}^\bullet, \tau)$ et un carré cartésien comme dans la
Définition \ref{definition-relative-dualizing-proper-flat}.
Nous interpréterons l'existence de l'isomorphisme
$(L(g')^*\omega_{X/Y}^\bullet, Lg^*\tau) \cong
(a'(\mathcal{O}_{Y'}), \text{Tr}_{f', \mathcal{O}_{Y'}})$
comme suit : elle affirme que, pour tout $M' \in D_\QCoh(\mathcal{O}_{X'})$,
le morphisme
$$
\Hom_{X'}(M', L(g')^*\omega_{X/Y}^\bullet)
\longrightarrow
\Hom_{Y'}(Rf'_*M', \mathcal{O}_{Y'}),\quad
\varphi' \longmapsto Lg^*\tau \circ Rf'_*\varphi'
$$
est un isomorphisme. Cela résulte de la définition de $a'$
et de la discussion de la section \ref{section-trace}. En particulier,
le lemme de Yoneda garantit l'unicité de l'isomorphisme.
\item Si $Y$ est lui-même affine, alors un complexe dualisant relatif
$(\omega_{X/Y}^\bullet, \tau)$ existe et est canoniquement isomorphe
à $(a(\mathcal{O}_Y), \text{Tr}_{f, \mathcal{O}_Y})$, où
$a$ est l'adjoint à droite de $Rf_*$ comme dans le
Lemme \ref{lemma-twisted-inverse-image}
et $\text{Tr}_f$ est comme dans la section \ref{section-trace}.
En effet, un diagramme comme dans la définition donne
un isomorphisme $L(g')^*a(\mathcal{O}_Y) \to a'(\mathcal{O}_{Y'})$ d'après le
Lemme \ref{lemma-more-base-change},
compatible aux morphismes trace d'après le
Lemme \ref{lemma-trace-map-and-base-change}.
\end{enumerate}
Cela fournit exactement les données nécessaires pour recoller les
complexes dualisants relatifs donnés localement en un complexe global. Nous suggérons au lecteur
de sauter les démonstrations des lemmes suivants.

\begin{lemma}
\label{lemma-relative-dualizing-RHom}
Soit $S$ un schéma. Soit $X \to Y$ un morphisme propre et plat
d'espaces algébriques, de présentation finie.
Si $(\omega_{X/Y}^\bullet, \tau)$ est un complexe dualisant relatif,
alors $\mathcal{O}_X \to
R\SheafHom_{\mathcal{O}_X}(\omega_{X/Y}^\bullet, \omega_{X/Y}^\bullet)$
est un isomorphisme et les faisceaux de cohomologie de $Rf_*\omega_{X/Y}^\bullet$
sont nuls en degrés positifs.
\end{lemma}

\begin{proof}
Il suffit de le démontrer après changement de base par un schéma affine \'etale
sur $Y$ ; cela résulte alors du
Lemme \ref{lemma-properties-relative-dualizing}.
\end{proof}

\begin{lemma}
\label{lemma-uniqueness-relative-dualizing}
Soit $S$ un schéma. Soit $X \to Y$ un morphisme propre et plat
d'espaces algébriques, de présentation finie.
Si $(\omega_j^\bullet, \tau_j)$, $j = 1, 2$,
sont deux complexes dualisants relatifs sur $X/Y$,
alors il existe un unique isomorphisme
$(\omega_1^\bullet, \tau_1) \to (\omega_2^\bullet, \tau_2)$.
\end{lemma}

\begin{proof}
Considérons $g : Y' \to Y$ \'etale, avec $Y'$ un schéma affine,
et notons $X' = Y' \times_Y X$ le changement de base.
D'après la Définition \ref{definition-relative-dualizing-proper-flat}
et la discussion qui la suit, il existe un unique isomorphisme
$\iota : (\omega_1^\bullet|_{X'}, \tau_1|_{Y'}) \to
(\omega_2^\bullet|_{X'}, \tau_2|_{Y'})$. Si $Y'' \to Y'$
est un autre morphisme \'etale d'affines et $X'' = Y'' \times_Y X$,
alors $\iota|_{X''}$ est l'unique isomorphisme
$(\omega_1^\bullet|_{X''}, \tau_1|_{Y''}) \to
(\omega_2^\bullet|_{X''}, \tau_2|_{Y''})$ (par unicité).
De plus, nous avons
$$
\text{Ext}^p_{X'}(\omega_1^\bullet|_{X'}, \omega_2^\bullet|_{X'}) = 0,
\quad p < 0
$$
car
$\mathcal{O}_{X'} \cong
R\SheafHom_{\mathcal{O}_{X'}}(\omega_1^\bullet|_{X'}, \omega_1^\bullet|_{X'})
\cong
R\SheafHom_{\mathcal{O}_{X'}}(\omega_1^\bullet|_{X'}, \omega_2^\bullet|_{X'})$
d'après le Lemme \ref{lemma-relative-dualizing-RHom}.

\medskip\noindent
Choisissons un hyperrecouvrement \'etale $b : V \to Y$ tel que chaque
$V_n = \coprod_{i \in I_n} Y_{n, i}$, avec $Y_{n, i}$ affine.
C'est possible d'après Hyperrecouvrements, Lemme
\ref{hypercovering-lemma-hypercovering-object} et
Remarque \ref{hypercovering-remark-take-unions-hypercovering-X}
(pour remplacer l'hyperrecouvrement produit
dans le lemme par celui qui possède des réunions disjointes en chaque degré).
Notons $X_{n, i} = Y_{n, i} \times_Y X$ et $U_n = V_n \times_Y X$,
de sorte que nous obtenons un hyperrecouvrement \'etale
$a : U \to X$ (Hyperrecouvrements, Lemme
\ref{hypercovering-lemma-hypercovering-morphism-sites})
avec $U_n = \coprod X_{n, i}$.
Les hypothèses d'Espaces simpliciaux, Lemme
\ref{spaces-simplicial-lemma-fppf-neg-ext-zero-hom},
sont satisfaites pour $a : U \to X$ et les complexes
$\omega_1^\bullet$ et $\omega_2^\bullet$.
Nous obtenons donc un unique morphisme
$\iota : \omega_1^\bullet \to \omega_2^\bullet$
dont la restriction à $X_{0, i}$ est l'unique
isomorphisme $(\omega_1^\bullet|_{X_{0, i}}, \tau_1|_{Y_{0, i}}) \to
(\omega_2^\bullet|_{X_{0, i}}, \tau_2|_{Y_{0, i}})$.
Il reste à vérifier que le diagramme
$$
\xymatrix{
Rf_*\omega_1^\bullet \ar[rd]_{\tau_1} \ar[rr]_{Rf_*\iota} & &
Rf_*\omega_1^\bullet \ar[ld]^{\tau_2} \\
& \mathcal{O}_Y
}
$$
est commutatif. Or nous savons que $Rf_*\omega_1^\bullet$ et
$Rf_*\omega_2^\bullet$ ont leurs faisceaux de cohomologie nuls en degrés
positifs (Lemme \ref{lemma-relative-dualizing-RHom}) ;
cette commutativité peut donc se vérifier après
restriction aux affines $Y_{0, i}$, où elle vaut par construction.
\end{proof}

\begin{lemma}
\label{lemma-covering-enough}
Soit $S$ un schéma. Soit $X \to Y$ un morphisme propre et plat
d'espaces algébriques, de présentation finie.
Soit $(\omega^\bullet, \tau)$ un couple formé
d'un objet $Y$-parfait de $D(\mathcal{O}_X)$ et d'un morphisme
$\tau : Rf_*\omega^\bullet \to \mathcal{O}_Y$.
Supposons donnés des diagrammes cartésiens
$$
\xymatrix{
X_i \ar[r]_{g_i'} \ar[d]_{f_i} & X \ar[d]^f \\
Y_i \ar[r]^{g_i} & Y
}
$$
avec $Y_i$ affine, tels que $\{g_i : Y_i \to Y\}$ soit un recouvrement \'etale,
et des isomorphismes de couples $(\omega^\bullet|_{X_i}, \tau|_{Y_i})
\to (a_i(\mathcal{O}_{Y_i}), \text{Tr}_{f_i, \mathcal{O}_{Y_i}})$
comme dans la Définition \ref{definition-relative-dualizing-proper-flat}.
Alors $(\omega^\bullet, \tau)$ est un complexe dualisant relatif pour $X$ sur $Y$.
\end{lemma}

\begin{proof}
Soient $g : Y' \to Y$ et $X', f', g', a'$ comme dans la
Définition \ref{definition-relative-dualizing-proper-flat}.
Posons $((\omega')^\bullet, \tau') = (L(g')^*\omega^\bullet, Lg^*\tau)$.
Nous pouvons trouver un recouvrement \'etale fini
$\{Y'_j \to Y'\}$ par des affines qui raffine $\{Y_i \times_Y Y' \to Y'\}$
(Topologies, Lemme \ref{topologies-lemma-etale-affine}).
Ainsi, pour chaque $j$, il existe un $i_j$ et un morphisme
$k_j : Y'_j \to Y_{i_j}$ sur $Y$. Considérons les produits fibrés
$$
\xymatrix{
X'_j \ar[r]_{h_j'} \ar[d]_{f'_j} &
X' \ar[d]^{f'} \\
Y'_j \ar[r]^{h_j} & Y'
}
$$
Notons $k'_j : X'_j \to X_{i_j}$ le morphisme induit (changement de base
de $k_j$ par $f_{i_j}$). En restreignant les isomorphismes donnés à $Y'_j$
par le morphisme $k'_j$, nous obtenons des isomorphismes de couples
$((\omega')^\bullet|_{X'_j}, \tau'|_{Y'_j})
\to (a_j(\mathcal{O}_{Y'_j}), \text{Tr}_{f'_j, \mathcal{O}_{Y'_j}})$.
Après avoir remplacé $f : X \to Y$ par $f' : X' \to Y'$, nous nous ramenons au
problème résolu dans le paragraphe suivant.

\medskip\noindent
Supposons $Y$ affine. Problème : montrer que $(\omega^\bullet, \tau)$ est
isomorphe à $(\omega_{X/Y}^\bullet, \text{Tr}) =
(a(\mathcal{O}_Y), \text{Tr}_{f, \mathcal{O}_Y})$.
Nous pouvons supposer que notre recouvrement $\{Y_i \to Y\}$ est donné par un unique
morphisme \'etale surjectif $\{g : Y' \to Y\}$ entre schémas affines.
En effet, nous pouvons d'abord remplacer $\{g_i: Y_i \to Y\}$ par un sous-recouvrement
fini, puis poser $g = \coprod g_i :  Y' = \coprod Y_i \to Y$ ;
nous omettons quelques détails. Posons $X' = Y' \times_Y X$, avec les morphismes
$f', g'$ de la Définition \ref{definition-relative-dualizing-proper-flat}.
La seule donnée qui subsiste est alors un isomorphisme
$$
(\omega^\bullet|_{X'}, \tau|_{Y'}) \to
(a'(\mathcal{O}_{Y'}), \text{Tr}_{f', \mathcal{O}_{Y'}})
$$
Puisque $(\omega_{X/Y}^\bullet, \text{Tr})$ est un complexe dualisant relatif
(voir la discussion qui suit la
Définition \ref{definition-relative-dualizing-proper-flat}),
il existe un unique isomorphisme
$$
(\omega_{X/Y}^\bullet|_{X'}, \text{Tr}|_{Y'}) \to
(a'(\mathcal{O}_{Y'}), \text{Tr}_{f', \mathcal{O}_{Y'}})
$$
L'unicité résulte par exemple du Lemme \ref{lemma-uniqueness-relative-dualizing}.
En combinant les isomorphismes affichés, nous trouvons un isomorphisme
$$
\alpha :
(\omega^\bullet|_{X'}, \tau|_{Y'}) \to
(\omega_{X/Y}^\bullet|_{X'}, \text{Tr}|_{Y'})
$$
Posons $Y'' = Y' \times_Y Y'$ et $X'' = Y'' \times_Y X$ ; les deux
images réciproques de $\alpha$ sur $X''$ doivent être égales, encore par unicité.
Comme les auto-Ext négatifs de
$\omega_{X'/Y'}^\bullet$ sur $X'$ sont nuls (Lemme \ref{lemma-relative-dualizing-RHom})
et que cette propriété subsiste après image réciproque par toute projection
$Y' \times_Y \ldots \times_Y Y' \to Y'$ (nous omettons un petit détail -- comparer
avec la démonstration du Lemme \ref{lemma-uniqueness-relative-dualizing}),
nous trouvons que $\alpha$ descend en un isomorphisme
$\omega^\bullet \to \omega_{X/Y}^\bullet$
sur $X$ d'après Espaces simpliciaux, Lemme
\ref{spaces-simplicial-lemma-fppf-neg-ext-zero-hom}.
\end{proof}

\begin{lemma}
\label{lemma-existence-relative-dualizing}
Soit $S$ un schéma. Soit $X \to Y$ un morphisme propre et plat
d'espaces algébriques, de présentation finie.
Il existe un complexe dualisant relatif $(\omega_{X/Y}^\bullet, \tau)$.
\end{lemma}

\begin{proof}
Choisissons un hyperrecouvrement \'etale $b : V \to Y$ tel que chaque
$V_n = \coprod_{i \in I_n} Y_{n, i}$, avec $Y_{n, i}$ affine.
C'est possible d'après Hyperrecouvrements, Lemme
\ref{hypercovering-lemma-hypercovering-object} et
Remarque \ref{hypercovering-remark-take-unions-hypercovering-X}
(pour remplacer l'hyperrecouvrement produit
dans le lemme par celui qui possède des réunions disjointes en chaque degré).
Notons $X_{n, i} = Y_{n, i} \times_Y X$ et $U_n = V_n \times_Y X$,
de sorte que nous obtenons un hyperrecouvrement \'etale
$a : U \to X$ (Hyperrecouvrements, Lemme
\ref{hypercovering-lemma-hypercovering-morphism-sites})
avec $U_n = \coprod X_{n, i}$.
Pour chaque $n, i$, il existe un complexe dualisant relatif
$(\omega_{n, i}^\bullet, \tau_{n, i})$ sur $X_{n, i}/Y_{n, i}$.
Voir la discussion qui suit la
Définition \ref{definition-relative-dualizing-proper-flat}.
Pour $\varphi : [m] \to [n]$ et $i \in I_n$, considérons les morphismes
$g_{\varphi, i} : Y_{n, i} \to Y_{m, \alpha(\varphi)}$ et
$g'_{\varphi, i} : X_{n, i} \to X_{m, \alpha(\varphi)}$
qui font partie de la structure des hyperrecouvrements donnés
(Hyperrecouvrements, section \ref{hypercovering-section-hypercovering-sites}).
Nous avons alors un unique isomorphisme
$$
\iota_{n, i, \varphi} :
(L(g'_{n, i})^*\omega_{n, i}^\bullet, Lg_{n, i}^*\tau_{n, i})
\longrightarrow
(\omega_{m, \alpha(\varphi)(i)}^\bullet, \tau_{m, \alpha(\varphi)(i)})
$$
de couples ; voir la discussion qui suit la
Définition \ref{definition-relative-dualizing-proper-flat}.
Observons que les auto-Ext négatifs de $\omega_{n, i}^\bullet$
sur $X_{n, i}$ sont nuls d'après le Lemme \ref{lemma-relative-dualizing-RHom}.
Notons $(\omega_n^\bullet, \tau_n)$ le couple sur $U_n/V_n$ construit à l'aide
des couples $(\omega_{n, i}^\bullet, \tau_{n, i})$ pour $i \in I_n$.
Pour $\varphi : [m] \to [n]$ et $i \in I_n$, considérons les morphismes
$g_\varphi : V_n \to V_m$ et $g'_\varphi : U_n \to U_m$ qui font partie
de la structure des espaces algébriques simpliciaux $V$ et $U$.
Nous avons alors des isomorphismes uniques
$$
\iota_\varphi :
(L(g'_\varphi)^*\omega_n^\bullet, Lg_\varphi^*\tau_n)
\longrightarrow
(\omega_m^\bullet, \tau_m)
$$
de couples, construits à partir des isomorphismes sur les morceaux.
L'unicité garantit que ces isomorphismes satisfont
à la condition de transitivité formulée dans
Espaces simpliciaux, Définition
\ref{spaces-simplicial-definition-cartesian-derived-modules}.
Les hypothèses d'Espaces simpliciaux, Lemme
\ref{spaces-simplicial-lemma-fppf-glue-neg-ext-zero},
sont satisfaites pour $a : U \to X$, les complexes $\omega_n^\bullet$
et les isomorphismes $\iota_\varphi$\footnote{Ce lemme utilise seulement
$\omega_0^\bullet$ et les deux morphismes $\delta_1^1, \delta_0^1 : [1] \to [0]$.
Le lecteur peut sauter les premières lignes de la démonstration du
lemme cité,
car nous disposons ici déjà d'un système simplicial de la
catégorie dérivée des modules.}.
Nous obtenons ainsi un objet $\omega^\bullet$ de $D_\QCoh(\mathcal{O}_X)$
muni d'un isomorphisme
$\iota_0 : \omega^\bullet|_{U_0} \to \omega_0^\bullet$
compatible aux deux isomorphismes
$\iota_{\delta^1_0}$ et $\iota_{\delta^1_1}$.
Enfin, nous appliquons
Espaces simpliciaux, Lemme
\ref{spaces-simplicial-lemma-fppf-neg-ext-zero-hom},
pour obtenir un unique morphisme
$$
\tau : Rf_*\omega^\bullet \longrightarrow \mathcal{O}_Y
$$
dont la restriction à $V_0$ coïncide avec $\tau_0$ ; nous omettons quelques détails --
comparer par exemple avec la fin de la démonstration du
Lemme \ref{lemma-uniqueness-relative-dualizing},
pour voir pourquoi nous avons
l'annulation requise des Ext négatifs.
D'après le Lemme \ref{lemma-covering-enough},
le couple $(\omega^\bullet, \tau)$ est un complexe dualisant relatif,
ce qui achève la démonstration.
\end{proof}

\begin{lemma}
\label{lemma-base-change-relative-dualizing}
Soit $S$ un schéma. Considérons un carré cartésien
$$
\xymatrix{
X' \ar[d]_{f'} \ar[r]_{g'} & X \ar[d]^f \\
Y' \ar[r]^g & Y
}
$$
d'espaces algébriques sur $S$. Supposons que $X \to Y$ soit propre, plat et
de présentation finie. Soit $(\omega_{X/Y}^\bullet, \tau)$ un
complexe dualisant relatif pour $f$. Alors
$(L(g')^*\omega_{X/Y}^\bullet, Lg^*\tau)$ est un complexe dualisant
relatif pour $f'$.
\end{lemma}

\begin{proof}
Observons que $L(g')^*\omega_{X/Y}^\bullet$ est $Y'$-parfait d'après
Compléments sur les morphismes d'espaces, Lemme
\ref{spaces-more-morphisms-lemma-base-change-relatively-perfect}.
L'autre condition de la
Définition \ref{definition-relative-dualizing-proper-flat}
résulte de la transitivité des produits fibrés.
\end{proof}






 
\section{Comparaison avec le cas des schémas}
\label{section-comparison}

\noindent
Il faudrait ajouter beaucoup plus de choses dans cette section.

\begin{lemma}
\label{lemma-compare}
Soit $S$ un schéma. Soit $f : X \to Y$ un morphisme d'espaces algébriques
quasi-compacts et quasi-séparés sur $S$.
Supposons $X$ et $Y$ représentables, et soit $f_0 : X_0 \to Y_0$ un
morphisme de schémas représentant $f$ (notation peu élégante, mais temporaire).
Soit $a : D_\QCoh(\mathcal{O}_Y) \to D_\QCoh(\mathcal{O}_X)$
l'adjoint à droite de $Rf_*$ du Lemme \ref{lemma-twisted-inverse-image}.
Soit $a_0 : D_\QCoh(\mathcal{O}_{Y_0}) \to D_\QCoh(\mathcal{O}_{X_0})$
l'adjoint à droite de $Rf_*$ de
Dualité pour les schémas, Lemme \ref{duality-lemma-twisted-inverse-image}.
Alors
$$
\xymatrix{
D_\QCoh(\mathcal{O}_{X_0})
\ar@{=}[rrrrrr]_{\text{Catégories dérivées des espaces, Lemme
\ref{spaces-perfect-lemma-derived-quasi-coherent-small-etale-site}}}
& & & & & &
D_\QCoh(\mathcal{O}_X) \\
D_\QCoh(\mathcal{O}_{Y_0}) \ar[u]^{a_0}
\ar@{=}[rrrrrr]^{\text{Catégories dérivées des espaces, Lemme
\ref{spaces-perfect-lemma-derived-quasi-coherent-small-etale-site}}}
& & & & & &
D_\QCoh(\mathcal{O}_Y) \ar[u]_a
}
$$
est commutatif.
\end{lemma}

\begin{proof}
Cela résulte de l'unicité des adjoints et des compatibilités de
Catégories dérivées des espaces, Remarque
\ref{spaces-perfect-remark-match-total-direct-images}.
\end{proof}









\begin{multicols}{2}[\section{Autres chapitres}]
\noindent
Préliminaires
\begin{enumerate}
\item \hyperref[introduction-section-phantom]{Introduction}
\item \hyperref[conventions-section-phantom]{Conventions}
\item \hyperref[sets-section-phantom]{Théorie des ensembles}
\item \hyperref[categories-section-phantom]{Catégories}
\item \hyperref[topology-section-phantom]{Topologie}
\item \hyperref[sheaves-section-phantom]{Faisceaux sur les espaces}
\item \hyperref[sites-section-phantom]{Sites et faisceaux}
\item \hyperref[stacks-section-phantom]{Champs}
\item \hyperref[fields-section-phantom]{Corps}
\item \hyperref[algebra-section-phantom]{Algèbre commutative}
\item \hyperref[brauer-section-phantom]{Groupes de Brauer}
\item \hyperref[homology-section-phantom]{Algèbre homologique}
\item \hyperref[derived-section-phantom]{Catégories dérivées}
\item \hyperref[simplicial-section-phantom]{Méthodes simpliciales}
\item \hyperref[more-algebra-section-phantom]{Compléments d'algèbre}
\item \hyperref[smoothing-section-phantom]{Lissage des morphismes d'anneaux}
\item \hyperref[modules-section-phantom]{Faisceaux de modules}
\item \hyperref[sites-modules-section-phantom]{Modules sur les sites}
\item \hyperref[injectives-section-phantom]{Objets injectifs}
\item \hyperref[cohomology-section-phantom]{Cohomologie des faisceaux}
\item \hyperref[sites-cohomology-section-phantom]{Cohomologie sur les sites}
\item \hyperref[dga-section-phantom]{Algèbre différentielle graduée}
\item \hyperref[dpa-section-phantom]{Algèbre à puissances divisées}
\item \hyperref[sdga-section-phantom]{Faisceaux différentiels gradués}
\item \hyperref[hypercovering-section-phantom]{Hyperrecouvrements}
\end{enumerate}
Schémas
\begin{enumerate}
\setcounter{enumi}{25}
\item \hyperref[schemes-section-phantom]{Schémas}
\item \hyperref[constructions-section-phantom]{Constructions de schémas}
\item \hyperref[properties-section-phantom]{Propriétés des schémas}
\item \hyperref[morphisms-section-phantom]{Morphismes de schémas}
\item \hyperref[coherent-section-phantom]{Cohomologie des schémas}
\item \hyperref[divisors-section-phantom]{Diviseurs}
\item \hyperref[limits-section-phantom]{Limites de schémas}
\item \hyperref[varieties-section-phantom]{Variétés}
\item \hyperref[topologies-section-phantom]{Topologies sur les schémas}
\item \hyperref[descent-section-phantom]{Descente}
\item \hyperref[perfect-section-phantom]{Catégories dérivées de schémas}
\item \hyperref[more-morphisms-section-phantom]{Compléments sur les morphismes}
\item \hyperref[flat-section-phantom]{Compléments sur la platitude}
\item \hyperref[groupoids-section-phantom]{Schémas en groupoïdes}
\item \hyperref[more-groupoids-section-phantom]{Compléments sur les schémas en groupoïdes}
\item \hyperref[etale-section-phantom]{Morphismes étales de schémas}
\end{enumerate}
Thèmes de théorie des schémas
\begin{enumerate}
\setcounter{enumi}{41}
\item \hyperref[chow-section-phantom]{Homologie de Chow}
\item \hyperref[intersection-section-phantom]{Théorie de l'intersection}
\item \hyperref[pic-section-phantom]{Schémas de Picard des courbes}
\item \hyperref[weil-section-phantom]{Théories cohomologiques de Weil}
\item \hyperref[adequate-section-phantom]{Modules adéquats}
\item \hyperref[dualizing-section-phantom]{Complexes dualisants}
\item \hyperref[duality-section-phantom]{Dualité pour les schémas}
\item \hyperref[discriminant-section-phantom]{Discriminants et différentes}
\item \hyperref[derham-section-phantom]{Cohomologie de de Rham}
\item \hyperref[local-cohomology-section-phantom]{Cohomologie locale}
\item \hyperref[algebraization-section-phantom]{Géométrie algébrique et formelle}
\item \hyperref[curves-section-phantom]{Courbes algébriques}
\item \hyperref[resolve-section-phantom]{Résolution des surfaces}
\item \hyperref[models-section-phantom]{Réduction semi-stable}
\item \hyperref[functors-section-phantom]{Foncteurs et morphismes}
\item \hyperref[equiv-section-phantom]{Catégories dérivées de variétés}
\item \hyperref[pione-section-phantom]{Groupes fondamentaux de schémas}
\item \hyperref[etale-cohomology-section-phantom]{Cohomologie étale}
\item \hyperref[crystalline-section-phantom]{Cohomologie cristalline}
\item \hyperref[proetale-section-phantom]{Cohomologie pro-étale}
\item \hyperref[relative-cycles-section-phantom]{Cycles relatifs}
\item \hyperref[more-etale-section-phantom]{Compléments de cohomologie étale}
\item \hyperref[trace-section-phantom]{La formule des traces}
\end{enumerate}
Espaces algébriques
\begin{enumerate}
\setcounter{enumi}{64}
\item \hyperref[spaces-section-phantom]{Espaces algébriques}
\item \hyperref[spaces-properties-section-phantom]{Propriétés des espaces algébriques}
\item \hyperref[spaces-morphisms-section-phantom]{Morphismes d'espaces algébriques}
\item \hyperref[decent-spaces-section-phantom]{Espaces algébriques décents}
\item \hyperref[spaces-cohomology-section-phantom]{Cohomologie des espaces algébriques}
\item \hyperref[spaces-limits-section-phantom]{Limites d'espaces algébriques}
\item \hyperref[spaces-divisors-section-phantom]{Diviseurs sur les espaces algébriques}
\item \hyperref[spaces-over-fields-section-phantom]{Espaces algébriques sur des corps}
\item \hyperref[spaces-topologies-section-phantom]{Topologies sur les espaces algébriques}
\item \hyperref[spaces-descent-section-phantom]{Descente et espaces algébriques}
\item \hyperref[spaces-perfect-section-phantom]{Catégories dérivées d'espaces}
\item \hyperref[spaces-more-morphisms-section-phantom]{Compléments sur les morphismes d'espaces}
\item \hyperref[spaces-flat-section-phantom]{Platitude sur les espaces algébriques}
\item \hyperref[spaces-groupoids-section-phantom]{Groupoïdes dans les espaces algébriques}
\item \hyperref[spaces-more-groupoids-section-phantom]{Compléments sur les groupoïdes dans les espaces}
\item \hyperref[bootstrap-section-phantom]{Amorçage}
\item \hyperref[spaces-pushouts-section-phantom]{Sommes amalgamées d'espaces algébriques}
\end{enumerate}
Thèmes de géométrie
\begin{enumerate}
\setcounter{enumi}{81}
\item \hyperref[spaces-chow-section-phantom]{Groupes de Chow des espaces}
\item \hyperref[groupoids-quotients-section-phantom]{Quotients de groupoïdes}
\item \hyperref[spaces-more-cohomology-section-phantom]{Compléments sur la cohomologie des espaces}
\item \hyperref[spaces-simplicial-section-phantom]{Espaces simpliciaux}
\item \hyperref[spaces-duality-section-phantom]{Dualité pour les espaces}
\item \hyperref[formal-spaces-section-phantom]{Espaces algébriques formels}
\item \hyperref[restricted-section-phantom]{Algébrisation des espaces formels}
\item \hyperref[spaces-resolve-section-phantom]{Résolution des surfaces revisitée}
\end{enumerate}
Théorie des déformations
\begin{enumerate}
\setcounter{enumi}{89}
\item \hyperref[formal-defos-section-phantom]{Théorie formelle des déformations}
\item \hyperref[defos-section-phantom]{Théorie des déformations}
\item \hyperref[cotangent-section-phantom]{Le complexe cotangent}
\item \hyperref[examples-defos-section-phantom]{Problèmes de déformation}
\end{enumerate}
Champs algébriques
\begin{enumerate}
\setcounter{enumi}{93}
\item \hyperref[algebraic-section-phantom]{Champs algébriques}
\item \hyperref[examples-stacks-section-phantom]{Exemples de champs}
\item \hyperref[stacks-sheaves-section-phantom]{Faisceaux sur les champs algébriques}
\item \hyperref[criteria-section-phantom]{Critères de représentabilité}
\item \hyperref[artin-section-phantom]{Axiomes d'Artin}
\item \hyperref[quot-section-phantom]{Espaces de Quot et de Hilbert}
\item \hyperref[stacks-properties-section-phantom]{Propriétés des champs algébriques}
\item \hyperref[stacks-morphisms-section-phantom]{Morphismes de champs algébriques}
\item \hyperref[stacks-limits-section-phantom]{Limites de champs algébriques}
\item \hyperref[stacks-cohomology-section-phantom]{Cohomologie des champs algébriques}
\item \hyperref[stacks-perfect-section-phantom]{Catégories dérivées de champs}
\item \hyperref[stacks-introduction-section-phantom]{Introduction aux champs algébriques}
\item \hyperref[stacks-more-morphisms-section-phantom]{Compléments sur les morphismes de champs}
\item \hyperref[stacks-geometry-section-phantom]{La géométrie des champs}
\end{enumerate}
Thèmes de théorie des espaces de modules
\begin{enumerate}
\setcounter{enumi}{107}
\item \hyperref[moduli-section-phantom]{Champs de modules}
\item \hyperref[moduli-curves-section-phantom]{Espaces de modules de courbes}
\end{enumerate}
Divers
\begin{enumerate}
\setcounter{enumi}{109}
\item \hyperref[examples-section-phantom]{Exemples}
\item \hyperref[exercises-section-phantom]{Exercices}
\item \hyperref[guide-section-phantom]{Guide bibliographique}
\item \hyperref[desirables-section-phantom]{Desiderata}
\item \hyperref[coding-section-phantom]{Style de programmation}
\item \hyperref[obsolete-section-phantom]{Obsolète}
\item \hyperref[fdl-section-phantom]{Licence de documentation libre GNU}
\item \hyperref[index-section-phantom]{Index généré automatiquement}
\end{enumerate}
\end{multicols}


\bibliography{my}
\bibliographystyle{amsalpha}

\end{document}
